% =============================================================================
% APPENDICES — Systematic Literature Review: Supplementary Materials
% Companion to Chapter~\ref{ch:slr}
%
% Usage note: include this file AFTER \appendix in the main thesis file:
%   \appendix
%   % =============================================================================
% APPENDICES — Systematic Literature Review: Supplementary Materials
% Companion to Chapter~\ref{ch:slr}
%
% Usage note: include this file AFTER \appendix in the main thesis file:
%   \appendix
%   % =============================================================================
% APPENDICES — Systematic Literature Review: Supplementary Materials
% Companion to Chapter~\ref{ch:slr}
%
% Usage note: include this file AFTER \appendix in the main thesis file:
%   \appendix
%   % =============================================================================
% APPENDICES — Systematic Literature Review: Supplementary Materials
% Companion to Chapter~\ref{ch:slr}
%
% Usage note: include this file AFTER \appendix in the main thesis file:
%   \appendix
%   \include{appendix_slr}
% The \appendix command causes subsequent \chapter{} calls to be labelled
% as APPENDIX A, B, C, etc. per ABNT NBR 14724.
% =============================================================================


% =============================================================================
% APPENDIX A — Search Protocol: Search Strings
% =============================================================================
\chapter{Search Protocol: Search Strings}
\label{app:A}

This appendix presents the four search strings executed across the five
digital libraries during the search phase of the PATHCAST SLR. Each string
was adapted to the syntax and field conventions of the target database; the
block structure (P and I) is preserved across all adaptations. Strings S1--S3
correspond to the principal, complementary (MSR), and frontier (stochastic
$\wedge$ PM) blocks, respectively; S4 is a supplementary string executed in
Scopus only, targeting Markov chain models in software testing.

\section{Search String Summary by Database}
\label{app:slr-strings}

\begin{table}[htbp]
\centering
\caption{Search strings by database and field scope.}
\label{tab:app-strings}
\small
\setlength{\tabcolsep}{4pt}
\begin{tabular}{p{2.4cm} c p{9.0cm}}
\toprule
\textbf{Database} & \textbf{Strings} & \textbf{Field / syntax notes} \\
\midrule
Scopus       & S1, S2, S3, S4 & \texttt{TITLE-ABS-KEY(...)};
                                  \texttt{DOCTYPE(ar OR cp)};
                                  \texttt{LANGUAGE(english)};
                                  1994--2026 \\
IEEE Xplore  & S1 (split)     & All Metadata; Boolean AND of sub-blocks;
                                  Conference + Journals; 1994--2026 \\
Web of Science & S1, S2, S3   & Topic \texttt{TS=(...)}; Article, Proceedings;
                                  English; 1994--2026 \\
SpringerLink & S1 (simplified)& Title + Abstract; 8 separate exports;
                                  Chapter + Article; English \\
ACM DL       & S1, S2         & Abstract; 2 queries; date 1994--2026 \\
\bottomrule
\end{tabular}
\end{table}

\section{String S1 — Principal ($\text{P} \wedge \text{I}$)}
\label{app:slr-s1}

Block P (software engineering context):

\begin{quote}
\small\ttfamily
("software engineering" OR "software development" OR "software process"
OR "software lifecycle" OR "software life cycle" OR "SDLC" OR "development
lifecycle" OR "development life cycle" OR "DevOps" OR "CI/CD" OR
"continuous integration" OR "continuous delivery" OR "agile development"
OR "agile process" OR "issue tracking" OR "bug tracking" OR "version
control" OR "code review" OR "pull request" OR "software project"
OR "software repository" OR "software maintenance" OR "software
evolution")
\end{quote}

Block I (techniques):

\begin{quote}
\small\ttfamily
("process mining" OR "process discovery" OR "conformance checking" OR
"workflow mining" OR "predictive process monitoring" OR "process prediction"
OR "remaining time prediction" OR "process forecasting" OR "event log" OR
"process simulation" OR "Monte Carlo simulation" OR "Markov chain" OR
"Markov model" OR "stochastic process model" OR "transition probability"
OR "absorbing chain" OR "lead time prediction" OR "cycle time prediction")
\end{quote}

\section{String S2 — Complementary (MSR $\rightarrow$ Process Modeling)}
\label{app:slr-s2}

\begin{quote}
\small\ttfamily
("mining software repositories" OR "software repository mining" OR
"GitHub mining" OR "Jira mining") AND ("process model" OR "Markov" OR
"Monte Carlo" OR "stochastic" OR "lead time" OR "cycle time" OR "throughput"
OR "transition matrix")
\end{quote}

\section{String S3 — Frontier (Stochastic $\wedge$ PM $\wedge$ SE)}
\label{app:slr-s3}

\begin{quote}
\small\ttfamily
("Monte Carlo" OR "Markov chain" OR "stochastic simulation" OR
"probabilistic forecast" OR "transition matrix") AND ("process mining" OR
"event log" OR "workflow mining" OR "process model" OR "business process")
AND ("software development" OR "software engineering" OR "software process"
OR "DevOps" OR "SDLC" OR "issue tracker" OR "commit log" OR "build pipeline")
\end{quote}

\section{String S4 — Markov Chain Software Testing (Scopus supplement)}
\label{app:slr-s4}

\begin{quote}
\small\ttfamily
("Markov chain" OR "Markov model" OR "hidden Markov" OR "stochastic automata")
AND ("software testing" OR "test coverage" OR "software reliability" OR
"defect prediction" OR "fault prediction")
\end{quote}


% =============================================================================
% APPENDIX B — LLM Screening Prompts
% =============================================================================
\chapter{LLM Screening Prompts}
\label{app:B}

The review employed \textit{claude-haiku-4-5-20251001} as primary screener in
two phases: title/abstract (T/A) and full-text (FT). Each phase uses a
\emph{system prompt} that encodes the scope, criteria, and output schema,
paired with a per-paper \emph{user prompt} that supplies the bibliographic
data. All prompts are version-controlled in
\texttt{config/screening\_criteria.py} in the review repository.

\section{Phase 1: Title/Abstract Screening}
\label{app:slr-ta-prompt}

\subsection*{System Prompt (T/A Phase)}

\begin{quote}
\small\ttfamily
You are an expert in systematic literature review (SLR) in the area of
process mining and software engineering. Your task is to perform the
title/abstract (T/A) screening for SLR PATHCAST.

\textbf{SLR PATHCAST scope:}
Intersection of (1) process mining, (2) stochastic modelling,
(3) forecasting/prediction, and (4) software development processes.
Focus: studies that analyse, discover, or predict software process
characteristics using event logs, repositories, or stochastic models.

\textbf{INCLUSION CRITERIA} (at least ONE must be satisfied to include):
IC1 -- Process Mining in Software: applies PM (discovery, conformance,
workflow mining, event log analysis, PPM) to SW artefacts (commits,
issues, PRs, CI/CD).
IC2 -- Stochastic Modelling of SW Processes: uses Markov chains, Monte
Carlo, stochastic Petri nets, or transition matrices to model SW dev
processes.
IC3 -- Process Metric Forecasting: predicts lead time, cycle time,
remaining time, throughput, or defect rates in SW processes using
data or event logs.
IC4 -- Repository Mining for Process: mines repositories (GitHub, Jira,
VCS) to discover or improve SW development process models.

\textbf{EXCLUSION CRITERIA} (any ONE is sufficient to exclude):
EC1 -- Non-SW domain: applied to healthcare, manufacturing, finance,
etc., with no link to SW processes.
EC2 -- SW as tool: "software" is the implementation platform, not the
process under study.
EC3 -- Algorithm without SW application: purely theoretical method,
no evaluation in SW processes.
EC4 -- Out-of-scope secondary study: survey/SLR not specifically
addressing PM or stochastic methods in SE.

\textbf{UNCERTAIN CASES (maybe):}
Relevant title without abstract; generic BPM article with possible SW
application; ambiguous terminology.

\textbf{STRUCTURED FIELDS} (use only the allowed enums):
evidence\_tags, software\_context, stochastic\_method,
forecast\_target, process\_data\_source, confidence.

Respond ONLY with valid JSON, no additional text.
\end{quote}

\subsection*{User Prompt Template (T/A Phase)}

\begin{quote}
\small\ttfamily
Evaluate the paper below for SLR PATHCAST.

\textbf{Title:} \{title\}
\textbf{Abstract:} \{abstract\}
\textbf{Venue/Type:} \{venue\} | \{doc\_type\} | \{year\}
\textbf{Source database:} \{source\_db\}
\textbf{Abstract origin:} \{abstract\_source\}

Respond with JSON in the exact format:
\{
  "decision": "include" | "exclude" | "maybe",
  "rationale": "<1-2 sentences justifying the decision>",
  "matched\_ic": ["IC1", "IC2"],
  "matched\_ec": ["EC1"],
  "evidence\_tags": ["process\_mining", "forecasting"],
  "software\_context": "software\_development\_process",
  "stochastic\_method": "markov\_chain",
  "forecast\_target": "remaining\_time",
  "process\_data\_source": "event\_logs",
  "confidence": "medium"
\}

Rules:
- "include" if $\geq$1 IC met and no decisive EC
- "exclude" if an EC clearly applies, or paper clearly meets no IC
- "maybe" if genuine doubt (title without abstract, domain ambiguity)
- matched\_ic and matched\_ec must be [] when not applicable
\end{quote}

\section{Phase 2: Full-Text Screening}
\label{app:slr-ft-prompt}

\subsection*{System Prompt (FT Phase)}

The FT system prompt is identical to the T/A system prompt above, with
three additional instructions that encode the stricter standard of the
full-text phase:

\begin{quote}
\small\ttfamily
THIS IS THE FULL-TEXT PHASE -- apply criteria rigorously:
- If the abstract clearly confirms $\geq$1 IC with no decisive EC
  $\rightarrow$ "include"
- If the abstract clearly confirms no IC applies, or an EC applies
  $\rightarrow$ "exclude"
- Use "pending" ONLY if the abstract is genuinely insufficient for
  a decision (absent or extremely short) AND the title suggests
  marginal relevance

When in doubt between include and exclude, prefer exclude -- this
phase requires confidence.

YOU MUST NEVER invent, complete, or infer a missing abstract.
\end{quote}

\subsection*{User Prompt Template (FT Phase)}

The FT user prompt extends the T/A template with a \textbf{T/A context block}
that supplies the previous-phase decision, rationale, and matched ICs:

\begin{quote}
\small\ttfamily
Evaluate the paper below for SLR PATHCAST -- FULL-TEXT screening (phase 2).

\textbf{Title:} \{title\}
\textbf{Abstract:} \{abstract\}
\textbf{Venue/Type:} \{venue\} | \{doc\_type\} | \{year\}
\textbf{Source database:} \{source\_db\}
\textbf{Abstract origin:} \{abstract\_source\}

\textbf{T/A phase context:}
  - Previous decision: \{ta\_decision\}
  - Previous rationale: \{ta\_rationale\}
  - ICs identified in phase 1: \{ta\_matched\_ic\}

Re-evaluate rigorously for the full-text phase.
Respond with JSON in the exact format:
\{
  "decision": "include" | "exclude" | "pending",
  [... same fields as T/A phase ...]
\}
\end{quote}

\section{Phase 2b: Screen-Blanks Pass}
\label{app:slr-blanks-prompt}

A third pass targeted the 186 papers without a decision after both prior
phases (primarily Band~E papers with no available abstract). The screen-blanks
pass reuses the same FT system and user prompts, with the abstract field
populated by the placeholder ``Abstract not available'' and the T/A context
block showing \texttt{ta\_decision~=~maybe} and
\texttt{ta\_matched\_ic~=~none}. Papers for which the LLM returned
\texttt{pending} and whose title analysis suggested low relevance were
subsequently closed under EC5 (full text inaccessible; $n = 148$).

\section{Model Parameters per Screening Phase}
\label{app:slr-prompt-params}

Table~\ref{tab:app-prompt-params} summarises the model parameters used
across all three screening phases and the data extraction step.

\begin{table}[htbp]
\centering
\caption{LLM parameters used at each screening phase.}
\label{tab:app-prompt-params}
\small
\setlength{\tabcolsep}{4pt}
\begin{tabular}{l l r r l}
\toprule
\textbf{Phase} & \textbf{Model} & \textbf{Max tok.} &
  \textbf{Batch} & \textbf{API} \\
\midrule
T/A screening      & claude-haiku-4-5-20251001 & 512  & 500 & Anthropic Batches \\
FT screening       & claude-haiku-4-5-20251001 & 512  & 100 & Anthropic Batches \\
Screen-blanks pass & claude-haiku-4-5-20251001 & 512  & 186 & Anthropic Batches \\
Data extraction    & claude-haiku-4-5-20251001 & 1024 & 100 & Anthropic Batches \\
Synthesis (ch.~3)  & claude-sonnet-4-6         & 2048 & --- & Anthropic Messages \\
\bottomrule
\end{tabular}
\end{table}


% =============================================================================
% APPENDIX C — Data Extraction Form
% =============================================================================
\chapter{Data Extraction Form}
\label{app:C}

Data extraction for all 169 included studies was performed with
\textit{claude-haiku-4-5-20251001} via the Anthropic Batches API. Papers
with a locally available PDF had their text pre-extracted with
\texttt{pdfplumber} (first 8 pages, maximum 6,000 characters). Papers
without a PDF used the enriched abstract from Semantic Scholar.
The prompt below was submitted once per paper.

\section{Extraction Prompt}
\label{app:slr-extraction-prompt}

\begin{quote}
\small\ttfamily
You are a researcher performing data extraction for SLR PATHCAST on
Process Mining and Stochastic Modelling in Software Processes.

\textbf{Title:} \{title\}
\textbf{Authors:} \{authors\}
\textbf{Year:} \{year\} | \textbf{Venue:} \{venue\}
\textbf{Inclusion criteria met:} \{ics\}
\{content\_block\}

Return ONLY a valid JSON object with exactly these fields:
\{
  "research\_question": "main research question (1-2 sentences)",
  "study\_type": "case\_study | experiment | survey | tool |
                 theoretical | simulation | mixed",
  "research\_contribution": "discovery | conformance | enhancement |
       prediction | simulation | framework | hybrid",
  "pm\_technique": "alpha | inductive\_miner | heuristic |
       conformance\_checking | social\_network | declarative | other | none",
  "stochastic\_technique": "markov\_chain | stochastic\_petri\_net |
       monte\_carlo | system\_dynamics | other | none",
  "software\_artifact": "commits | issues | ci\_cd | ide\_logs |
       vcs | jira | github | other | none",
  "software\_process": "development | testing | maintenance |
       code\_review | bug\_fixing | deployment | requirements | other",
  "dataset\_source": "open\_source | industrial | academic |
                     synthetic | not\_specified",
  "dataset\_public": "yes | no | partial | not\_mentioned",
  "tool\_used": "tool names (e.g.\ ProM, pm4py, Disco) or not\_mentioned",
  "main\_finding": "main result/contribution (2-3 sentences)",
  "limitations": "reported limitations (1-2 sentences) or not\_mentioned",
  "replication\_package": "yes | no | partial | not\_mentioned"
\}

Respond ONLY with JSON, no additional text.
\end{quote}

\section{Extraction Form Fields}
\label{app:slr-extraction-fields}

Table~\ref{tab:app-extraction-fields} describes all fields of the structured
extraction form stored in \texttt{extraction\_template.csv}.

\begin{table}[htbp]
\centering
\caption{Data extraction form: all fields in \texttt{extraction\_template.csv}.}
\label{tab:app-extraction-fields}
\footnotesize
\setlength{\tabcolsep}{3pt}
\begin{tabular}{p{3.2cm} p{4.4cm} p{4.6cm}}
\toprule
\textbf{Field} & \textbf{Description} & \textbf{Allowed values / format} \\
\midrule
\multicolumn{3}{l}{\textit{Identification}} \\
\texttt{internal\_id}   & 8-char hash (dedup key)      & --- \\
\texttt{title}          & Paper title                  & free text \\
\texttt{doi}            & Digital Object Identifier    & URI \\
\texttt{year}           & Publication year             & YYYY \\
\texttt{source\_db}     & Source database              & scopus, ieee, acm, springer, wos, snowball \\
\texttt{ft\_matched\_ic}& Inclusion criteria met       & IC1$|$IC2$|$IC3$|$IC4 \\
\texttt{ft\_screened\_by}& Who made FT decision        & llm, manual, llm\_pdf \\
\midrule
\multicolumn{3}{l}{\textit{Enriched metadata (Semantic Scholar)}} \\
\texttt{authors}        & Author list                  & semicolon-separated \\
\texttt{venue}          & Publication venue            & free text \\
\texttt{journal\_name}  & Journal/conference name      & free text \\
\texttt{abstract}       & Paper abstract               & free text \\
\texttt{pdf\_available} & PDF locally present          & yes / no \\
\midrule
\multicolumn{3}{l}{\textit{LLM-extracted fields}} \\
\texttt{research\_question} & Main RQ               & 1--2 sentences \\
\texttt{study\_type}        & Study design          & case\_study $|$ experiment $|$ survey $|$ tool $|$ theoretical $|$ simulation $|$ mixed \\
\texttt{research\_contribution} & Contribution type & discovery $|$ conformance $|$ enhancement $|$ prediction $|$ simulation $|$ framework $|$ hybrid \\
\texttt{pm\_technique}      & PM algorithm used     & alpha $|$ inductive\_miner $|$ heuristic $|$ conformance\_checking $|$ social\_network $|$ declarative $|$ other $|$ none \\
\texttt{stochastic\_technique} & Stochastic method  & markov\_chain $|$ stochastic\_petri\_net $|$ monte\_carlo $|$ system\_dynamics $|$ other $|$ none \\
\texttt{software\_artifact} & Data source           & commits $|$ issues $|$ ci\_cd $|$ ide\_logs $|$ vcs $|$ jira $|$ github $|$ other $|$ none \\
\texttt{software\_process}  & SDLC phase            & development $|$ testing $|$ maintenance $|$ code\_review $|$ bug\_fixing $|$ deployment $|$ requirements $|$ other \\
\texttt{dataset\_source}    & Dataset provenance    & open\_source $|$ industrial $|$ academic $|$ synthetic $|$ not\_specified \\
\texttt{dataset\_public}    & Data availability     & yes $|$ no $|$ partial $|$ not\_mentioned \\
\texttt{tool\_used}         & Tools/platforms       & free text \\
\texttt{main\_finding}      & Key result            & 2--3 sentences \\
\texttt{limitations}        & Study limitations     & 1--2 sentences \\
\texttt{replication\_package} & Artefacts available & yes $|$ no $|$ partial $|$ not\_mentioned \\
\midrule
\multicolumn{3}{l}{\textit{Manual fields (reading phase)}} \\
\texttt{quality\_score}     & QA checklist score (0--8) & numeric \\
\texttt{extraction\_notes}  & Free-form reviewer note   & free text \\
\bottomrule
\end{tabular}
\end{table}

\section{Structured Output Controlled Vocabulary}
\label{app:slr-schema}

Table~\ref{tab:app-enum-schema} defines the complete controlled vocabulary
for the structured fields returned by the LLM at both screening phases.
These enums are version-controlled in \texttt{config/screening\_criteria.py}
and were not changed between screening rounds to ensure consistency.

\begin{table}[htbp]
\centering
\caption{Controlled vocabulary (enums) for LLM structured output.}
\label{tab:app-enum-schema}
\small
\setlength{\tabcolsep}{4pt}
\begin{tabular}{p{3.2cm} p{9.4cm}}
\toprule
\textbf{Field} & \textbf{Allowed values} \\
\midrule
\texttt{decision} (T/A) &
  \texttt{include}, \texttt{exclude}, \texttt{maybe} \\
\texttt{decision} (FT) &
  \texttt{include}, \texttt{exclude}, \texttt{pending} \\
\texttt{software\_context} &
  \texttt{software\_development\_process}, \texttt{repository\_mining},
  \texttt{ci\_cd}, \texttt{issue\_bug\_workflow}, \texttt{software\_testing},
  \texttt{requirements\_engineering}, \texttt{software\_project\_management},
  \texttt{unclear}, \texttt{not\_software\_process} \\
\texttt{stochastic\_method} &
  \texttt{none}, \texttt{markov\_chain}, \texttt{hidden\_markov\_model},
  \texttt{monte\_carlo}, \texttt{stochastic\_petri\_net},
  \texttt{bayesian\_model}, \texttt{probabilistic\_model},
  \texttt{simulation}, \texttt{queueing\_model},
  \texttt{other\_stochastic}, \texttt{unclear} \\
\texttt{forecast\_target} &
  \texttt{none}, \texttt{lead\_time}, \texttt{cycle\_time},
  \texttt{remaining\_time}, \texttt{throughput}, \texttt{defect\_rate},
  \texttt{build\_outcome}, \texttt{reliability},
  \texttt{completion\_time}, \texttt{other\_process\_metric},
  \texttt{unclear} \\
\texttt{process\_data\_source} &
  \texttt{none}, \texttt{event\_logs}, \texttt{version\_control},
  \texttt{issue\_tracker}, \texttt{pull\_requests}, \texttt{ci\_cd\_logs},
  \texttt{software\_repository\_mixed}, \texttt{synthetic\_data},
  \texttt{simulated\_process}, \texttt{survey\_or\_secondary},
  \texttt{unclear} \\
\texttt{confidence} &
  \texttt{low}, \texttt{medium}, \texttt{high} \\
\texttt{evidence\_tags} &
  \texttt{process\_mining}, \texttt{software\_process},
  \texttt{repository\_mining}, \texttt{stochastic\_modeling},
  \texttt{forecasting}, \texttt{event\_log}, \texttt{version\_control},
  \texttt{issue\_tracking}, \texttt{pull\_requests}, \texttt{ci\_cd},
  \texttt{markov}, \texttt{hidden\_markov\_model}, \texttt{monte\_carlo},
  \texttt{stochastic\_petri\_net}, \texttt{bayesian\_model},
  \texttt{simulation}, \texttt{lead\_time}, \texttt{cycle\_time},
  \texttt{remaining\_time}, \texttt{throughput}, \texttt{defect\_prediction},
  \texttt{build\_prediction}, \texttt{reliability},
  \texttt{insufficient\_abstract}, \texttt{full\_text\_required} \\
\bottomrule
\end{tabular}
\end{table}


% =============================================================================
% APPENDIX D — Assessment and Classification Instruments
% =============================================================================
\chapter{Assessment and Classification Instruments}
\label{app:D}

This appendix gathers the instruments used to classify, enrich, and evaluate
papers throughout the screening and selection process: the priority band
taxonomy, the abstract enrichment cascade, the quality assessment checklist,
the control paper set, and the SPMF integration-level taxonomy.

\section{Priority Bands for the Full-Text Review Queue}
\label{app:slr-bands}

Prior to full-text screening, the 886-paper queue was stratified into five
priority bands (A--E) to optimise resource allocation. Abstract enrichment
and LLM screening were applied in band order.
Table~\ref{tab:app-bands} defines the classification rules.

\begin{table}[htbp]
\centering
\caption{Priority band definitions for the full-text review queue.}
\label{tab:app-bands}
\small
\setlength{\tabcolsep}{4pt}
\begin{tabular}{c r p{8.5cm}}
\toprule
\textbf{Band} & \textbf{$n$} & \textbf{Classification rule} \\
\midrule
A &  95 & T/A decision = \texttt{include} and $\geq$2 distinct ICs matched \\
B &  20 & T/A decision = \texttt{include} and exactly 1 IC matched \\
C &  39 & T/A decision = \texttt{maybe} and abstract available \\
D & 152 & T/A decision = \texttt{maybe}, at least 1 IC signal in title, no abstract \\
E & 584 & T/A decision = \texttt{maybe}, no IC signal in title, no abstract \\
\midrule
\textbf{Total} & \textbf{886} & \\
\bottomrule
\end{tabular}
\end{table}

Band~D papers with a DOI were prioritised for manual PDF retrieval. Band~E
papers without a DOI and without an abstract were closed as EC5 (full text
inaccessible) at the end of the screen-blanks pass.

\section{Abstract Enrichment Cascade}
\label{app:slr-enrichment}

To increase the number of papers assessable by the LLM, a four-source
cascade was executed prior to full-text screening.
Algorithm~\ref{alg:enrichment} describes the procedure.

\begin{figure}[htbp]
\centering
\small
\begin{tabular}{l}
\hline
\textbf{Algorithm 1} \quad Abstract Enrichment Cascade \\
\hline
\textbf{Input:} Paper record $p$ with \texttt{doi} and/or \texttt{title} \\
\textbf{Output:} Updated $p$ with \texttt{abstract}, \texttt{authors}, \texttt{venue} \\
\hline
\textbf{1.} If $p$.\texttt{abstract} is non-empty $\rightarrow$ \textbf{return} $p$ \\
\textbf{2.} If $p$.\texttt{doi} is present: \\
\quad\textbf{2a.} Query \textbf{Semantic Scholar} Batch API (\texttt{DOI:\{doi\}}) \\
\quad\quad If match found: update \texttt{abstract}, \texttt{authors}, \texttt{venue}; \textbf{return} \\
\quad\textbf{2b.} Query \textbf{OpenAlex} (\texttt{filter=doi:\{doi\}}) \\
\quad\quad If match found: update \texttt{abstract}; \textbf{return} \\
\quad\textbf{2c.} Query \textbf{CORE} (\texttt{/search/works?q=doi:\{doi\}}) \\
\quad\quad If match found: update \texttt{abstract}; \textbf{return} \\
\quad\textbf{2d.} Query \textbf{Crossref} (\texttt{/works/\{doi\}}) \\
\quad\quad If abstract present: update; \textbf{return} \\
\textbf{3.} If $p$.\texttt{doi} is absent: \\
\quad\textbf{3a.} Query \textbf{Semantic Scholar} by title (fuzzy) \\
\quad\quad If similarity $\geq 0.85$: update; \textbf{return} \\
\textbf{4.} Mark $p$.\texttt{abstract\_source} = \texttt{missing\_or\_unverified} \\
\hline
\end{tabular}
\caption{Abstract enrichment cascade applied to all 886 full-text queue papers.}
\label{alg:enrichment}
\end{figure}

The cascade raised abstract coverage from 16.6\% (147/886) to 72.6\%
(643/886) before Round~2 screening. Of the 643 enriched abstracts,
590 were sourced from Semantic Scholar, 31 from OpenAlex, 13 from CORE,
and 9 from Crossref or the title-based fallback.
Table~\ref{tab:app-enrichment-stats} summarises the enrichment outcome
by priority band.

\begin{table}[htbp]
\centering
\caption{Abstract enrichment results by priority band.}
\label{tab:app-enrichment-stats}
\small
\setlength{\tabcolsep}{4pt}
\begin{tabular}{c r r r r}
\toprule
\textbf{Band} & \textbf{Papers} & \textbf{With abstract (pre)} &
  \textbf{Gained} & \textbf{Coverage (post)} \\
\midrule
A &  95 & 91 (95.8\%) &   4 &  95 (100.0\%) \\
B &  20 & 18 (90.0\%) &   2 &  20 (100.0\%) \\
C &  39 & 39 (100\%)  &   0 &  39 (100.0\%) \\
D & 152 &  0  (0.0\%) & 118 & 118  (77.6\%) \\
E & 584 &  0  (0.0\%) & 371 & 371  (63.5\%) \\
\midrule
\textbf{Total} & \textbf{886} & \textbf{148 (16.7\%)} &
  \textbf{495} & \textbf{643 (72.6\%)} \\
\bottomrule
\end{tabular}
\end{table}

\section{Quality Assessment Checklist}
\label{app:slr-qa}

The eight-item quality checklist (Table~\ref{tab:app-qa-full}) is applied
to each included primary study during the manual reading phase. Items
QA1--QA4 assess methodological rigour; QA5--QA7 assess empirical soundness
and reproducibility; QA8 specifically targets process model quality reporting,
which is central to PATHCAST's evaluation strategy. Studies scoring below
4/8 are excluded from the synthesis.

\begin{table}[htbp]
\centering
\caption{Quality assessment checklist with scoring guidance.}
\label{tab:app-qa-full}
\small
\setlength{\tabcolsep}{4pt}
\begin{tabular}{c p{5.5cm} p{5.5cm}}
\toprule
\textbf{ID} & \textbf{Criterion} & \textbf{Scoring guidance} \\
\midrule
QA1 & Are the research objectives clearly stated? &
  Score 1 if RQs or hypotheses are explicitly enumerated; 0 otherwise. \\
QA2 & Is the software engineering context described in sufficient detail? &
  Score 1 if the SDLC phase, artefact type, and process scope are defined. \\
QA3 & Is the data source (event log or repository) described reproducibly? &
  Score 1 if dataset origin, collection method, and key statistics are given. \\
QA4 & Is the PM or stochastic technique formally defined? &
  Score 1 if algorithm, parameters, and configuration are specified. \\
QA5 & Are results validated empirically? &
  Score 1 if validation uses real or semi-real data; 0 for simulation-only. \\
QA6 & Are threats to validity discussed? &
  Score 1 if at least internal and external validity are addressed. \\
QA7 & Is the study reproducible (data and/or code available)? &
  Score 1 if a replication package or open dataset is linked or described. \\
QA8 & Are process model quality metrics reported (fitness, precision)? &
  Score 1 if at least one formal quality metric is quantified. \\
\midrule
\multicolumn{2}{l}{\textbf{Total score range}} & 0--8 (threshold: $\geq$4) \\
\bottomrule
\end{tabular}
\end{table}

Current QA status: 76 of 169 studies assessed (full Round~1 set); mean
score = 5.18 (SD = 1.13), median = 5.0 (IQR: 4--6); 5 studies scored
below the threshold and were excluded from the synthesis. The remaining
93 studies (Rounds~2+) are pending manual QA assessment.

\section{Control Paper Set}
\label{app:slr-control}

The 12-paper control set was used to validate the search strings before
execution. Table~\ref{tab:app-control} lists all control papers with their
IC profile and the source through which each was captured.

\begin{table}[htbp]
\centering
\caption{Complete control paper set for search string validation.}
\label{tab:app-control}
\small
\setlength{\tabcolsep}{3pt}
\begin{tabular}{c p{4.0cm} p{3.8cm} l l}
\toprule
\textbf{\#} & \textbf{Paper} & \textbf{Justification} & \textbf{ICs} &
  \textbf{Captured via} \\
\midrule
V1  & Cook \& Wolf (1998)~\cite{cook1998}              & Seminal PM + SE             & IC1, IC2 & ACM \\
V2  & Rubin et al.\ (2007)~\cite{rubin2007}            & First named PM framework for SW & IC1 & Scopus \\
V3  & Poncin et al.\ (2011)~\cite{poncin2011}          & PM + MSR intersection       & IC1, IC4 & IEEE \\
V4  & Caldeira \& Cardoso (2019)                        & PM + SW process assessment  & IC1      & IEEE \\
V5  & Lemos et al.\ (2011)                              & PM for SW management        & IC1      & IEEE \\
V6  & Whittaker \& Thomason (1994)~\cite{whittaker1994} & Markov + SW testing (seminal) & IC2    & Scopus \\
V7  & Tax et al.\ (2017)~\cite{tax2017}                & PPM seminal (BPM baseline)  & IC3      & Control \\
V8  & Rogge-Solti \& Weske (2015)                       & Stochastic + PM forecasting & IC2, IC3 & Scopus \\
V9  & Rozinat et al.\ (2009)                            & PM $\rightarrow$ Simulation  & IC1, IC2 & Control \\
V10 & Gupta (2017)~\cite{gupta2017}                    & PM multi-repository SE      & IC1, IC3 & ACM \\
V11 & Jalote \& Kamma (2021)~\cite{jalote2021}         & IDE logs + Markov chains    & IC1, IC2 & Scopus \\
V12 & Ortu et al.\ (2023)~\cite{ortu2023}              & GitHub Markov + SNA         & IC2, IC4 & Scopus \\
\midrule
\multicolumn{4}{l}{Capture rate from primary queries alone} & 10/12 (83.3\%) \\
\multicolumn{4}{l}{Capture rate including snowballing}      & 12/12 (100\%) \\
\bottomrule
\end{tabular}
\end{table}

Papers V4, V5, V8, and V9 were not in the initial working set but were
recovered through snowballing during the FT phase; they are included in
the final corpus.

\section{SPMF Integration-Level Taxonomy}
\label{app:slr-spmf}

The \textit{Software Process Mining and Forecasting} (SPMF) taxonomy, used
throughout Chapter~\ref{ch:slr}, classifies studies along three orthogonal
dimensions (Table~\ref{tab:app-spmf}).

\begin{table}[htbp]
\centering
\caption{SPMF taxonomy: three dimensions with level definitions.}
\label{tab:app-spmf}
\small
\setlength{\tabcolsep}{4pt}
\begin{tabular}{l l p{7.8cm}}
\toprule
\textbf{Dimension} & \textbf{Level} & \textbf{Definition} \\
\midrule
\multirow{4}{*}{\textbf{Analytical Depth}}
  & AD1 -- Descriptive  & Discovery of as-is process models; conformance checking \\
  & AD2 -- Diagnostic   & Bottleneck detection; rework quantification; root-cause analysis \\
  & AD3 -- Predictive   & Outcome prediction; remaining-time estimation \\
  & AD4 -- Forecasting  & Distributional forecasts; probabilistic completion estimates \\
\midrule
\multirow{5}{*}{\textbf{SDLC Coverage}}
  & SC-R & Requirements and design \\
  & SC-D & Development (coding, commits, pull requests) \\
  & SC-Q & Quality (testing, reviews, defects) \\
  & SC-C & CI/CD and delivery \\
  & SC-O & Operations and maintenance \\
\midrule
\multirow{4}{*}{\textbf{Technique Integration}}
  & TI-L0 & Single technique (PM only, ML only, or stochastic only) \\
  & TI-L1 & Two techniques present in same study without explicit pipeline contract \\
  & TI-L2 & Partial bridge: two techniques linked, but incomplete or missing ML closure \\
  & TI-L3 & Unified pipeline: PM $\rightarrow$ Markov $\rightarrow$ MC $\rightarrow$ ML
             with formal inter-stage contracts \\
\bottomrule
\end{tabular}
\end{table}

The L0--L3 scale in dimension TI maps to the integration-level coding used
in Sections~\ref{sec:slr-rq3} and~\ref{app:slr-top30}. PATHCAST is positioned
at TI-L3, AD4, with SC coverage spanning SC-D, SC-Q, and SC-C.


% =============================================================================
% APPENDIX E — Reading List and Pipeline Artefacts
% =============================================================================
\chapter{Reading List and Pipeline Artefacts}
\label{app:E}

This appendix presents the ranked reading list of 32 primary studies
selected for in-depth synthesis and documents the scripts and persistent
artefacts of the review pipeline.

\section{Multi-Criteria Scoring Rubric}
\label{app:slr-ranking-criteria}

Table~\ref{tab:app-ranking-criteria} presents the 15-point rubric used to
rank studies. Three seminal papers (Cook 1995, Cook 1998, Rubin 2007) are
included regardless of score.

\begin{table}[htbp]
\centering
\caption{Multi-criteria scoring rubric for the top-32 reading list.}
\label{tab:app-ranking-criteria}
\small
\setlength{\tabcolsep}{4pt}
\begin{tabular}{p{7.5cm} c}
\toprule
\textbf{Criterion} & \textbf{Points} \\
\midrule
IC2 present (stochastic modelling in SE)          & +3 \\
IC3 present (forecasting of process metrics)      & +3 \\
IC1 present (process mining in SE)                & +2 \\
IC4 present (repository mining for process)       & +1 \\
PDF locally available                             & +2 \\
Published in 2020 or later                        & +1 \\
High-impact venue (IEEE/ACM/IST/JSS/IS/TSE)      & +1 \\
Empirical study type (case\_study or experiment)  & +1 \\
Contribution type = prediction or simulation      & +1 \\
\midrule
\textbf{Maximum score}                            & \textbf{15} \\
\bottomrule
\end{tabular}
\end{table}

\section{Top-32 Ranked Reading List}
\label{app:slr-top30}

\begin{table}[p]
\centering
\caption{Top-32 reading list with score, IC cluster, and main finding
  summary. Papers marked $\dagger$ are forced-inclusion seminal works.}
\label{tab:app-top30}
\small
\setlength{\tabcolsep}{3pt}
\begin{tabular}{c c p{3.8cm} c c p{5.7cm}}
\toprule
\textbf{Rank} & \textbf{Pts} & \textbf{Authors (year)} &
  \textbf{ICs} & \textbf{PDF} & \textbf{Main finding (summary)} \\
\midrule
 1 & 14 & Joshi \& Kahani (2024)~\cite{joshi2024}
         & IC2,IC3,IC4 & \checkmark
         & RL/MDP for GitHub PR outcome prediction \\
 2 & 14 & Li et al.\ (2020)~\cite{li2020}
         & IC2,IC3 & \checkmark
         & Hybrid SD+DES simulation for industrial project forecasting \\
 3 & 13 & Guinea-Cabrera et al.\ (2025)~\cite{guinea2025}
         & IC1,IC2,IC3 & \checkmark
         & PRIMAD Digital Twin for SDLC (pm4py + Akka + LightGBM), L2 \\
 4 & 13 & Ortu et al.\ (2023)~\cite{ortu2023}
         & IC2,IC4 & \checkmark
         & Fault-insertion/fixing Markov + SNA in Apache projects \\
 5 & 13 & Ben Mesmia et al.\ (2021)~\cite{benmesmia2021}
         & IC2,IC3 & \checkmark
         & Non-Markovian SPN for DevOps workflow duration prediction \\
 6 & 12 & Buliga et al.\ (2025)~\cite{buliga2025}
         & IC1,IC3 & \checkmark
         & Generative AI + business process simulation for what-if \\
 7 & 12 & Bhadra et al.\ (2023)~\cite{bhadra2023}
         & IC2,IC3 & \checkmark
         & SPN model of CI/CD with reliability sensitivity analysis \\
 8 & 12 & Bhadra (2022)~\cite{bhadra2022}
         & IC2,IC3 & \checkmark
         & SPN model of CI/CD pipeline delivery probability \\
 9 & 12 & Massitela et al.\ (2018)~\cite{massitela2018}
         & IC2,IC3 & \checkmark
         & Stochastic Automata Network for Waterfall team estimation \\
10 & 11 & Incerto et al.\ (2025)~\cite{incerto2025}
         & IC1,IC2 & \checkmark
         & VLMC stochastic conformance; outperforms 18 state-of-art techniques \\
11 & 11 & Jo \& Kwon (2024)~\cite{jo2024}
         & IC2,IC4 & \checkmark
         & Probabilistic model checking (CTL) on GitHub repos \\
12 & 11 & Jo et al.\ (2023)~\cite{jo2023}
         & IC2,IC4 & \checkmark
         & DTMC collaboration patterns in GitHub (5 repos) \\
13 & 11 & L\'opez-Pintado \& Dumas (2023)~\cite{lopezpintado2023}
         & IC1,IC2 & ---
         & Probabilistic resource calendars discovered from event logs \\
14 & 11 & Caldeira et al.\ (2022)~\cite{caldeira2022}
         & IC1,IC3 & \checkmark
         & Process complexity $\to$ cyclomatic complexity ($r \approx 0.43$) \\
15 & 11 & Jalote \& Kamma (2021)~\cite{jalote2021}
         & IC1,IC2 & \checkmark
         & Markov chains on IDE logs for programmer productivity \\
16 & 11 & Lunesu et al.\ (2021)~\cite{lunesu2021}
         & IC2,IC3 & \checkmark
         & Monte Carlo risk assessment for agile software development \\
17 & 11 & Nafreen et al.\ (2020)~\cite{nafreen2020}
         & IC2,IC3 & ---
         & SRGM + defect tracking for hybrid reliability forecasting \\
18 & 11 & Gupta (2017)~\cite{gupta2017}
         & IC1,IC3,IC4 & ---
         & Inductive miner + predictive analytics for SW maintenance \\
19 & 10 & Salmani et al.\ (2022)
         & IC1,IC3 & \checkmark
         & Process mining for requirements engineering in MDPs \\
20 & 10 & Jeon et al.\ (2015)~\cite{jeon2015}
         & IC2,IC3,IC4 & ---
         & Markov-based probabilistic risk prediction from repo data \\
21 & 10 & Tian \& Noore (2005)~\cite{tian2005}
         & IC2,IC3 & ---
         & Markov defect-removal modelling under code churn \\
22 & 10 & Cook \& Wolf (1995)$^\dagger$~\cite{cook1995}
         & IC1,IC2 & \checkmark
         & First automated process discovery from event traces \\
23 &  9 & Pourbafrani et al.\ (2021)~\cite{pourbafrani2021}
         & IC1,IC3 & \checkmark
         & SIMPT: process tree + interactive what-if simulation tool \\
24 &  9 & Tyagi et al.\ (2021)~\cite{tyagi2021}
         & IC2,IC3 & ---
         & State-space Markov framework for agile planning and tracking \\
25 &  9 & Razi et al.\ (2019)
         & IC1,IC2 & \checkmark
         & Conformance checking for SDLC process compliance \\
26 &  9 & Krismayer et al.\ (2019)
         & IC1,IC4 & \checkmark
         & Declarative PM for repository-based process monitoring \\
27 &  9 & Baskoro et al.\ (2019)
         & IC3,IC4 & ---
         & ML-based effort prediction from repository data \\
28 &  9 & (Anon.\ 2017)
         & IC2,IC3 & ---
         & Probabilistic defect prediction for software product lines \\
29 &  9 & Magennis (2015)~\cite{magennis2015}
         & IC2,IC3 & ---
         & Monte Carlo cycle-time economic analysis \\
30 &  9 & Washizaki et al.\ (2015)~\cite{washizaki2015}
         & IC2,IC3 & ---
         & GSRM release-time prediction for open-source projects \\
31 &  9 & Cook \& Wolf (1998)$^\dagger$~\cite{cook1998}
         & IC1,IC2 & ---
         & Markov/FSA/NN process discovery from event-based data \\
32 &  2 & Rubin et al.\ (2007)$^\dagger$~\cite{rubin2007}
         & IC1 & \checkmark
         & First named PM framework for SDLC (ProM + SCM logs) \\
\bottomrule
\end{tabular}
\end{table}

\section{Review Pipeline Scripts}
\label{app:slr-pipeline}

The review was conducted using a custom Python pipeline (\texttt{pipeline/})
operating on a structured CSV data store. Table~\ref{tab:app-pipeline}
documents each script with its function and main invocation command.

\begin{table}[htbp]
\centering
\caption{Review pipeline scripts and their roles.}
\label{tab:app-pipeline}
\small
\setlength{\tabcolsep}{3pt}
\begin{tabular}{p{4.2cm} p{4.5cm} p{4.0cm}}
\toprule
\textbf{Script} & \textbf{Function} & \textbf{Main command} \\
\midrule
\texttt{pipeline/dedup.py} &
  Cross-source deduplication (DOI + fuzzy title) &
  \texttt{python main.py dedup} \\[2pt]
\texttt{pipeline/enrich.py} &
  Abstract enrichment cascade (4 sources) &
  \texttt{python main.py enrich} \\[2pt]
\texttt{pipeline/screening.py} &
  T/A LLM screening via Batches API &
  \texttt{python main.py screen -{}-run} \\[2pt]
\texttt{pipeline/fulltext.py} &
  FT LLM screening + manual review &
  \texttt{python main.py fulltext -{}-run} \\[2pt]
\texttt{pipeline/pdf\_downloader.py} &
  PDF retrieval (Unpaywall, S2, OA) &
  \texttt{python main.py pdf} \\[2pt]
\texttt{pipeline/extract\_prep.py} &
  S2 metadata enrichment + PDF copy &
  \texttt{python pipeline/} \texttt{extract\_prep.py} \\[2pt]
\texttt{pipeline/extract\_llm.py} &
  Structured data extraction (LLM, 169 papers) &
  \texttt{python pipeline/} \texttt{extract\_llm.py -{}-run} \\[2pt]
\texttt{pipeline/synth\_llm.py} &
  Synthesis paragraph generation per IC cluster &
  \texttt{python pipeline/} \texttt{synth\_llm.py} \\[2pt]
\texttt{pipeline/finalization.py} &
  PRISMA snapshot, EC5 closure, artefact export &
  \texttt{python main.py finalize} \\
\bottomrule
\end{tabular}
\end{table}

\section{Persistent Artefact Registry}
\label{app:slr-artefacts}

Table~\ref{tab:app-artefacts} lists the persistent artefacts produced by the
review, with their role, location, and current state.

\begin{table}[htbp]
\centering
\caption{Persistent artefacts of the review.}
\label{tab:app-artefacts}
\small
\setlength{\tabcolsep}{3pt}
\begin{tabular}{p{6.5cm} p{3.0cm} p{2.8cm}}
\toprule
\textbf{File} & \textbf{Role} & \textbf{State} \\
\midrule
\texttt{results/screening/} \texttt{ft\_screening\_results.csv}
  & 886-paper decision log & Final \\[3pt]
\texttt{results/extraction/} \texttt{extraction\_template.csv}
  & 169-paper extraction data & 169/169 \\[3pt]
\texttt{results/final\_review/} \texttt{included\_studies\_current.csv}
  & Included study set & 169 rows \\[3pt]
\texttt{results/final\_review/} \texttt{top30\_reading\_list.csv}
  & Ranked reading list & 32 rows \\[3pt]
\texttt{results/final\_review/} \texttt{cap3\_synthesis\_v2.tex}
  & LLM synthesis paragraphs & For review \\[3pt]
\texttt{results/final\_review/} \texttt{missing\_references.bib}
  & BibTeX for 32 missing keys & Ready to merge \\[3pt]
\texttt{results/final\_review/} \texttt{pending\_inaccessible\_closed.csv}
  & EC5 closure record & 148 rows \\[3pt]
\texttt{results/final\_review/} \texttt{prisma\_summary\_current.txt}
  & PRISMA snapshot & Final \\[3pt]
\texttt{results/} \texttt{SLR\_EXECUCAO\_ATE\_AGORA.md}
  & Execution log (all sessions) & Up to date \\
\bottomrule
\end{tabular}
\end{table}

% The \appendix command causes subsequent \chapter{} calls to be labelled
% as APPENDIX A, B, C, etc. per ABNT NBR 14724.
% =============================================================================


% =============================================================================
% APPENDIX A — Search Protocol: Search Strings
% =============================================================================
\chapter{Search Protocol: Search Strings}
\label{app:A}

This appendix presents the four search strings executed across the five
digital libraries during the search phase of the PATHCAST SLR. Each string
was adapted to the syntax and field conventions of the target database; the
block structure (P and I) is preserved across all adaptations. Strings S1--S3
correspond to the principal, complementary (MSR), and frontier (stochastic
$\wedge$ PM) blocks, respectively; S4 is a supplementary string executed in
Scopus only, targeting Markov chain models in software testing.

\section{Search String Summary by Database}
\label{app:slr-strings}

\begin{table}[htbp]
\centering
\caption{Search strings by database and field scope.}
\label{tab:app-strings}
\small
\setlength{\tabcolsep}{4pt}
\begin{tabular}{p{2.4cm} c p{9.0cm}}
\toprule
\textbf{Database} & \textbf{Strings} & \textbf{Field / syntax notes} \\
\midrule
Scopus       & S1, S2, S3, S4 & \texttt{TITLE-ABS-KEY(...)};
                                  \texttt{DOCTYPE(ar OR cp)};
                                  \texttt{LANGUAGE(english)};
                                  1994--2026 \\
IEEE Xplore  & S1 (split)     & All Metadata; Boolean AND of sub-blocks;
                                  Conference + Journals; 1994--2026 \\
Web of Science & S1, S2, S3   & Topic \texttt{TS=(...)}; Article, Proceedings;
                                  English; 1994--2026 \\
SpringerLink & S1 (simplified)& Title + Abstract; 8 separate exports;
                                  Chapter + Article; English \\
ACM DL       & S1, S2         & Abstract; 2 queries; date 1994--2026 \\
\bottomrule
\end{tabular}
\end{table}

\section{String S1 — Principal ($\text{P} \wedge \text{I}$)}
\label{app:slr-s1}

Block P (software engineering context):

\begin{quote}
\small\ttfamily
("software engineering" OR "software development" OR "software process"
OR "software lifecycle" OR "software life cycle" OR "SDLC" OR "development
lifecycle" OR "development life cycle" OR "DevOps" OR "CI/CD" OR
"continuous integration" OR "continuous delivery" OR "agile development"
OR "agile process" OR "issue tracking" OR "bug tracking" OR "version
control" OR "code review" OR "pull request" OR "software project"
OR "software repository" OR "software maintenance" OR "software
evolution")
\end{quote}

Block I (techniques):

\begin{quote}
\small\ttfamily
("process mining" OR "process discovery" OR "conformance checking" OR
"workflow mining" OR "predictive process monitoring" OR "process prediction"
OR "remaining time prediction" OR "process forecasting" OR "event log" OR
"process simulation" OR "Monte Carlo simulation" OR "Markov chain" OR
"Markov model" OR "stochastic process model" OR "transition probability"
OR "absorbing chain" OR "lead time prediction" OR "cycle time prediction")
\end{quote}

\section{String S2 — Complementary (MSR $\rightarrow$ Process Modeling)}
\label{app:slr-s2}

\begin{quote}
\small\ttfamily
("mining software repositories" OR "software repository mining" OR
"GitHub mining" OR "Jira mining") AND ("process model" OR "Markov" OR
"Monte Carlo" OR "stochastic" OR "lead time" OR "cycle time" OR "throughput"
OR "transition matrix")
\end{quote}

\section{String S3 — Frontier (Stochastic $\wedge$ PM $\wedge$ SE)}
\label{app:slr-s3}

\begin{quote}
\small\ttfamily
("Monte Carlo" OR "Markov chain" OR "stochastic simulation" OR
"probabilistic forecast" OR "transition matrix") AND ("process mining" OR
"event log" OR "workflow mining" OR "process model" OR "business process")
AND ("software development" OR "software engineering" OR "software process"
OR "DevOps" OR "SDLC" OR "issue tracker" OR "commit log" OR "build pipeline")
\end{quote}

\section{String S4 — Markov Chain Software Testing (Scopus supplement)}
\label{app:slr-s4}

\begin{quote}
\small\ttfamily
("Markov chain" OR "Markov model" OR "hidden Markov" OR "stochastic automata")
AND ("software testing" OR "test coverage" OR "software reliability" OR
"defect prediction" OR "fault prediction")
\end{quote}


% =============================================================================
% APPENDIX B — LLM Screening Prompts
% =============================================================================
\chapter{LLM Screening Prompts}
\label{app:B}

The review employed \textit{claude-haiku-4-5-20251001} as primary screener in
two phases: title/abstract (T/A) and full-text (FT). Each phase uses a
\emph{system prompt} that encodes the scope, criteria, and output schema,
paired with a per-paper \emph{user prompt} that supplies the bibliographic
data. All prompts are version-controlled in
\texttt{config/screening\_criteria.py} in the review repository.

\section{Phase 1: Title/Abstract Screening}
\label{app:slr-ta-prompt}

\subsection*{System Prompt (T/A Phase)}

\begin{quote}
\small\ttfamily
You are an expert in systematic literature review (SLR) in the area of
process mining and software engineering. Your task is to perform the
title/abstract (T/A) screening for SLR PATHCAST.

\textbf{SLR PATHCAST scope:}
Intersection of (1) process mining, (2) stochastic modelling,
(3) forecasting/prediction, and (4) software development processes.
Focus: studies that analyse, discover, or predict software process
characteristics using event logs, repositories, or stochastic models.

\textbf{INCLUSION CRITERIA} (at least ONE must be satisfied to include):
IC1 -- Process Mining in Software: applies PM (discovery, conformance,
workflow mining, event log analysis, PPM) to SW artefacts (commits,
issues, PRs, CI/CD).
IC2 -- Stochastic Modelling of SW Processes: uses Markov chains, Monte
Carlo, stochastic Petri nets, or transition matrices to model SW dev
processes.
IC3 -- Process Metric Forecasting: predicts lead time, cycle time,
remaining time, throughput, or defect rates in SW processes using
data or event logs.
IC4 -- Repository Mining for Process: mines repositories (GitHub, Jira,
VCS) to discover or improve SW development process models.

\textbf{EXCLUSION CRITERIA} (any ONE is sufficient to exclude):
EC1 -- Non-SW domain: applied to healthcare, manufacturing, finance,
etc., with no link to SW processes.
EC2 -- SW as tool: "software" is the implementation platform, not the
process under study.
EC3 -- Algorithm without SW application: purely theoretical method,
no evaluation in SW processes.
EC4 -- Out-of-scope secondary study: survey/SLR not specifically
addressing PM or stochastic methods in SE.

\textbf{UNCERTAIN CASES (maybe):}
Relevant title without abstract; generic BPM article with possible SW
application; ambiguous terminology.

\textbf{STRUCTURED FIELDS} (use only the allowed enums):
evidence\_tags, software\_context, stochastic\_method,
forecast\_target, process\_data\_source, confidence.

Respond ONLY with valid JSON, no additional text.
\end{quote}

\subsection*{User Prompt Template (T/A Phase)}

\begin{quote}
\small\ttfamily
Evaluate the paper below for SLR PATHCAST.

\textbf{Title:} \{title\}
\textbf{Abstract:} \{abstract\}
\textbf{Venue/Type:} \{venue\} | \{doc\_type\} | \{year\}
\textbf{Source database:} \{source\_db\}
\textbf{Abstract origin:} \{abstract\_source\}

Respond with JSON in the exact format:
\{
  "decision": "include" | "exclude" | "maybe",
  "rationale": "<1-2 sentences justifying the decision>",
  "matched\_ic": ["IC1", "IC2"],
  "matched\_ec": ["EC1"],
  "evidence\_tags": ["process\_mining", "forecasting"],
  "software\_context": "software\_development\_process",
  "stochastic\_method": "markov\_chain",
  "forecast\_target": "remaining\_time",
  "process\_data\_source": "event\_logs",
  "confidence": "medium"
\}

Rules:
- "include" if $\geq$1 IC met and no decisive EC
- "exclude" if an EC clearly applies, or paper clearly meets no IC
- "maybe" if genuine doubt (title without abstract, domain ambiguity)
- matched\_ic and matched\_ec must be [] when not applicable
\end{quote}

\section{Phase 2: Full-Text Screening}
\label{app:slr-ft-prompt}

\subsection*{System Prompt (FT Phase)}

The FT system prompt is identical to the T/A system prompt above, with
three additional instructions that encode the stricter standard of the
full-text phase:

\begin{quote}
\small\ttfamily
THIS IS THE FULL-TEXT PHASE -- apply criteria rigorously:
- If the abstract clearly confirms $\geq$1 IC with no decisive EC
  $\rightarrow$ "include"
- If the abstract clearly confirms no IC applies, or an EC applies
  $\rightarrow$ "exclude"
- Use "pending" ONLY if the abstract is genuinely insufficient for
  a decision (absent or extremely short) AND the title suggests
  marginal relevance

When in doubt between include and exclude, prefer exclude -- this
phase requires confidence.

YOU MUST NEVER invent, complete, or infer a missing abstract.
\end{quote}

\subsection*{User Prompt Template (FT Phase)}

The FT user prompt extends the T/A template with a \textbf{T/A context block}
that supplies the previous-phase decision, rationale, and matched ICs:

\begin{quote}
\small\ttfamily
Evaluate the paper below for SLR PATHCAST -- FULL-TEXT screening (phase 2).

\textbf{Title:} \{title\}
\textbf{Abstract:} \{abstract\}
\textbf{Venue/Type:} \{venue\} | \{doc\_type\} | \{year\}
\textbf{Source database:} \{source\_db\}
\textbf{Abstract origin:} \{abstract\_source\}

\textbf{T/A phase context:}
  - Previous decision: \{ta\_decision\}
  - Previous rationale: \{ta\_rationale\}
  - ICs identified in phase 1: \{ta\_matched\_ic\}

Re-evaluate rigorously for the full-text phase.
Respond with JSON in the exact format:
\{
  "decision": "include" | "exclude" | "pending",
  [... same fields as T/A phase ...]
\}
\end{quote}

\section{Phase 2b: Screen-Blanks Pass}
\label{app:slr-blanks-prompt}

A third pass targeted the 186 papers without a decision after both prior
phases (primarily Band~E papers with no available abstract). The screen-blanks
pass reuses the same FT system and user prompts, with the abstract field
populated by the placeholder ``Abstract not available'' and the T/A context
block showing \texttt{ta\_decision~=~maybe} and
\texttt{ta\_matched\_ic~=~none}. Papers for which the LLM returned
\texttt{pending} and whose title analysis suggested low relevance were
subsequently closed under EC5 (full text inaccessible; $n = 148$).

\section{Model Parameters per Screening Phase}
\label{app:slr-prompt-params}

Table~\ref{tab:app-prompt-params} summarises the model parameters used
across all three screening phases and the data extraction step.

\begin{table}[htbp]
\centering
\caption{LLM parameters used at each screening phase.}
\label{tab:app-prompt-params}
\small
\setlength{\tabcolsep}{4pt}
\begin{tabular}{l l r r l}
\toprule
\textbf{Phase} & \textbf{Model} & \textbf{Max tok.} &
  \textbf{Batch} & \textbf{API} \\
\midrule
T/A screening      & claude-haiku-4-5-20251001 & 512  & 500 & Anthropic Batches \\
FT screening       & claude-haiku-4-5-20251001 & 512  & 100 & Anthropic Batches \\
Screen-blanks pass & claude-haiku-4-5-20251001 & 512  & 186 & Anthropic Batches \\
Data extraction    & claude-haiku-4-5-20251001 & 1024 & 100 & Anthropic Batches \\
Synthesis (ch.~3)  & claude-sonnet-4-6         & 2048 & --- & Anthropic Messages \\
\bottomrule
\end{tabular}
\end{table}


% =============================================================================
% APPENDIX C — Data Extraction Form
% =============================================================================
\chapter{Data Extraction Form}
\label{app:C}

Data extraction for all 169 included studies was performed with
\textit{claude-haiku-4-5-20251001} via the Anthropic Batches API. Papers
with a locally available PDF had their text pre-extracted with
\texttt{pdfplumber} (first 8 pages, maximum 6,000 characters). Papers
without a PDF used the enriched abstract from Semantic Scholar.
The prompt below was submitted once per paper.

\section{Extraction Prompt}
\label{app:slr-extraction-prompt}

\begin{quote}
\small\ttfamily
You are a researcher performing data extraction for SLR PATHCAST on
Process Mining and Stochastic Modelling in Software Processes.

\textbf{Title:} \{title\}
\textbf{Authors:} \{authors\}
\textbf{Year:} \{year\} | \textbf{Venue:} \{venue\}
\textbf{Inclusion criteria met:} \{ics\}
\{content\_block\}

Return ONLY a valid JSON object with exactly these fields:
\{
  "research\_question": "main research question (1-2 sentences)",
  "study\_type": "case\_study | experiment | survey | tool |
                 theoretical | simulation | mixed",
  "research\_contribution": "discovery | conformance | enhancement |
       prediction | simulation | framework | hybrid",
  "pm\_technique": "alpha | inductive\_miner | heuristic |
       conformance\_checking | social\_network | declarative | other | none",
  "stochastic\_technique": "markov\_chain | stochastic\_petri\_net |
       monte\_carlo | system\_dynamics | other | none",
  "software\_artifact": "commits | issues | ci\_cd | ide\_logs |
       vcs | jira | github | other | none",
  "software\_process": "development | testing | maintenance |
       code\_review | bug\_fixing | deployment | requirements | other",
  "dataset\_source": "open\_source | industrial | academic |
                     synthetic | not\_specified",
  "dataset\_public": "yes | no | partial | not\_mentioned",
  "tool\_used": "tool names (e.g.\ ProM, pm4py, Disco) or not\_mentioned",
  "main\_finding": "main result/contribution (2-3 sentences)",
  "limitations": "reported limitations (1-2 sentences) or not\_mentioned",
  "replication\_package": "yes | no | partial | not\_mentioned"
\}

Respond ONLY with JSON, no additional text.
\end{quote}

\section{Extraction Form Fields}
\label{app:slr-extraction-fields}

Table~\ref{tab:app-extraction-fields} describes all fields of the structured
extraction form stored in \texttt{extraction\_template.csv}.

\begin{table}[htbp]
\centering
\caption{Data extraction form: all fields in \texttt{extraction\_template.csv}.}
\label{tab:app-extraction-fields}
\footnotesize
\setlength{\tabcolsep}{3pt}
\begin{tabular}{p{3.2cm} p{4.4cm} p{4.6cm}}
\toprule
\textbf{Field} & \textbf{Description} & \textbf{Allowed values / format} \\
\midrule
\multicolumn{3}{l}{\textit{Identification}} \\
\texttt{internal\_id}   & 8-char hash (dedup key)      & --- \\
\texttt{title}          & Paper title                  & free text \\
\texttt{doi}            & Digital Object Identifier    & URI \\
\texttt{year}           & Publication year             & YYYY \\
\texttt{source\_db}     & Source database              & scopus, ieee, acm, springer, wos, snowball \\
\texttt{ft\_matched\_ic}& Inclusion criteria met       & IC1$|$IC2$|$IC3$|$IC4 \\
\texttt{ft\_screened\_by}& Who made FT decision        & llm, manual, llm\_pdf \\
\midrule
\multicolumn{3}{l}{\textit{Enriched metadata (Semantic Scholar)}} \\
\texttt{authors}        & Author list                  & semicolon-separated \\
\texttt{venue}          & Publication venue            & free text \\
\texttt{journal\_name}  & Journal/conference name      & free text \\
\texttt{abstract}       & Paper abstract               & free text \\
\texttt{pdf\_available} & PDF locally present          & yes / no \\
\midrule
\multicolumn{3}{l}{\textit{LLM-extracted fields}} \\
\texttt{research\_question} & Main RQ               & 1--2 sentences \\
\texttt{study\_type}        & Study design          & case\_study $|$ experiment $|$ survey $|$ tool $|$ theoretical $|$ simulation $|$ mixed \\
\texttt{research\_contribution} & Contribution type & discovery $|$ conformance $|$ enhancement $|$ prediction $|$ simulation $|$ framework $|$ hybrid \\
\texttt{pm\_technique}      & PM algorithm used     & alpha $|$ inductive\_miner $|$ heuristic $|$ conformance\_checking $|$ social\_network $|$ declarative $|$ other $|$ none \\
\texttt{stochastic\_technique} & Stochastic method  & markov\_chain $|$ stochastic\_petri\_net $|$ monte\_carlo $|$ system\_dynamics $|$ other $|$ none \\
\texttt{software\_artifact} & Data source           & commits $|$ issues $|$ ci\_cd $|$ ide\_logs $|$ vcs $|$ jira $|$ github $|$ other $|$ none \\
\texttt{software\_process}  & SDLC phase            & development $|$ testing $|$ maintenance $|$ code\_review $|$ bug\_fixing $|$ deployment $|$ requirements $|$ other \\
\texttt{dataset\_source}    & Dataset provenance    & open\_source $|$ industrial $|$ academic $|$ synthetic $|$ not\_specified \\
\texttt{dataset\_public}    & Data availability     & yes $|$ no $|$ partial $|$ not\_mentioned \\
\texttt{tool\_used}         & Tools/platforms       & free text \\
\texttt{main\_finding}      & Key result            & 2--3 sentences \\
\texttt{limitations}        & Study limitations     & 1--2 sentences \\
\texttt{replication\_package} & Artefacts available & yes $|$ no $|$ partial $|$ not\_mentioned \\
\midrule
\multicolumn{3}{l}{\textit{Manual fields (reading phase)}} \\
\texttt{quality\_score}     & QA checklist score (0--8) & numeric \\
\texttt{extraction\_notes}  & Free-form reviewer note   & free text \\
\bottomrule
\end{tabular}
\end{table}

\section{Structured Output Controlled Vocabulary}
\label{app:slr-schema}

Table~\ref{tab:app-enum-schema} defines the complete controlled vocabulary
for the structured fields returned by the LLM at both screening phases.
These enums are version-controlled in \texttt{config/screening\_criteria.py}
and were not changed between screening rounds to ensure consistency.

\begin{table}[htbp]
\centering
\caption{Controlled vocabulary (enums) for LLM structured output.}
\label{tab:app-enum-schema}
\small
\setlength{\tabcolsep}{4pt}
\begin{tabular}{p{3.2cm} p{9.4cm}}
\toprule
\textbf{Field} & \textbf{Allowed values} \\
\midrule
\texttt{decision} (T/A) &
  \texttt{include}, \texttt{exclude}, \texttt{maybe} \\
\texttt{decision} (FT) &
  \texttt{include}, \texttt{exclude}, \texttt{pending} \\
\texttt{software\_context} &
  \texttt{software\_development\_process}, \texttt{repository\_mining},
  \texttt{ci\_cd}, \texttt{issue\_bug\_workflow}, \texttt{software\_testing},
  \texttt{requirements\_engineering}, \texttt{software\_project\_management},
  \texttt{unclear}, \texttt{not\_software\_process} \\
\texttt{stochastic\_method} &
  \texttt{none}, \texttt{markov\_chain}, \texttt{hidden\_markov\_model},
  \texttt{monte\_carlo}, \texttt{stochastic\_petri\_net},
  \texttt{bayesian\_model}, \texttt{probabilistic\_model},
  \texttt{simulation}, \texttt{queueing\_model},
  \texttt{other\_stochastic}, \texttt{unclear} \\
\texttt{forecast\_target} &
  \texttt{none}, \texttt{lead\_time}, \texttt{cycle\_time},
  \texttt{remaining\_time}, \texttt{throughput}, \texttt{defect\_rate},
  \texttt{build\_outcome}, \texttt{reliability},
  \texttt{completion\_time}, \texttt{other\_process\_metric},
  \texttt{unclear} \\
\texttt{process\_data\_source} &
  \texttt{none}, \texttt{event\_logs}, \texttt{version\_control},
  \texttt{issue\_tracker}, \texttt{pull\_requests}, \texttt{ci\_cd\_logs},
  \texttt{software\_repository\_mixed}, \texttt{synthetic\_data},
  \texttt{simulated\_process}, \texttt{survey\_or\_secondary},
  \texttt{unclear} \\
\texttt{confidence} &
  \texttt{low}, \texttt{medium}, \texttt{high} \\
\texttt{evidence\_tags} &
  \texttt{process\_mining}, \texttt{software\_process},
  \texttt{repository\_mining}, \texttt{stochastic\_modeling},
  \texttt{forecasting}, \texttt{event\_log}, \texttt{version\_control},
  \texttt{issue\_tracking}, \texttt{pull\_requests}, \texttt{ci\_cd},
  \texttt{markov}, \texttt{hidden\_markov\_model}, \texttt{monte\_carlo},
  \texttt{stochastic\_petri\_net}, \texttt{bayesian\_model},
  \texttt{simulation}, \texttt{lead\_time}, \texttt{cycle\_time},
  \texttt{remaining\_time}, \texttt{throughput}, \texttt{defect\_prediction},
  \texttt{build\_prediction}, \texttt{reliability},
  \texttt{insufficient\_abstract}, \texttt{full\_text\_required} \\
\bottomrule
\end{tabular}
\end{table}


% =============================================================================
% APPENDIX D — Assessment and Classification Instruments
% =============================================================================
\chapter{Assessment and Classification Instruments}
\label{app:D}

This appendix gathers the instruments used to classify, enrich, and evaluate
papers throughout the screening and selection process: the priority band
taxonomy, the abstract enrichment cascade, the quality assessment checklist,
the control paper set, and the SPMF integration-level taxonomy.

\section{Priority Bands for the Full-Text Review Queue}
\label{app:slr-bands}

Prior to full-text screening, the 886-paper queue was stratified into five
priority bands (A--E) to optimise resource allocation. Abstract enrichment
and LLM screening were applied in band order.
Table~\ref{tab:app-bands} defines the classification rules.

\begin{table}[htbp]
\centering
\caption{Priority band definitions for the full-text review queue.}
\label{tab:app-bands}
\small
\setlength{\tabcolsep}{4pt}
\begin{tabular}{c r p{8.5cm}}
\toprule
\textbf{Band} & \textbf{$n$} & \textbf{Classification rule} \\
\midrule
A &  95 & T/A decision = \texttt{include} and $\geq$2 distinct ICs matched \\
B &  20 & T/A decision = \texttt{include} and exactly 1 IC matched \\
C &  39 & T/A decision = \texttt{maybe} and abstract available \\
D & 152 & T/A decision = \texttt{maybe}, at least 1 IC signal in title, no abstract \\
E & 584 & T/A decision = \texttt{maybe}, no IC signal in title, no abstract \\
\midrule
\textbf{Total} & \textbf{886} & \\
\bottomrule
\end{tabular}
\end{table}

Band~D papers with a DOI were prioritised for manual PDF retrieval. Band~E
papers without a DOI and without an abstract were closed as EC5 (full text
inaccessible) at the end of the screen-blanks pass.

\section{Abstract Enrichment Cascade}
\label{app:slr-enrichment}

To increase the number of papers assessable by the LLM, a four-source
cascade was executed prior to full-text screening.
Algorithm~\ref{alg:enrichment} describes the procedure.

\begin{figure}[htbp]
\centering
\small
\begin{tabular}{l}
\hline
\textbf{Algorithm 1} \quad Abstract Enrichment Cascade \\
\hline
\textbf{Input:} Paper record $p$ with \texttt{doi} and/or \texttt{title} \\
\textbf{Output:} Updated $p$ with \texttt{abstract}, \texttt{authors}, \texttt{venue} \\
\hline
\textbf{1.} If $p$.\texttt{abstract} is non-empty $\rightarrow$ \textbf{return} $p$ \\
\textbf{2.} If $p$.\texttt{doi} is present: \\
\quad\textbf{2a.} Query \textbf{Semantic Scholar} Batch API (\texttt{DOI:\{doi\}}) \\
\quad\quad If match found: update \texttt{abstract}, \texttt{authors}, \texttt{venue}; \textbf{return} \\
\quad\textbf{2b.} Query \textbf{OpenAlex} (\texttt{filter=doi:\{doi\}}) \\
\quad\quad If match found: update \texttt{abstract}; \textbf{return} \\
\quad\textbf{2c.} Query \textbf{CORE} (\texttt{/search/works?q=doi:\{doi\}}) \\
\quad\quad If match found: update \texttt{abstract}; \textbf{return} \\
\quad\textbf{2d.} Query \textbf{Crossref} (\texttt{/works/\{doi\}}) \\
\quad\quad If abstract present: update; \textbf{return} \\
\textbf{3.} If $p$.\texttt{doi} is absent: \\
\quad\textbf{3a.} Query \textbf{Semantic Scholar} by title (fuzzy) \\
\quad\quad If similarity $\geq 0.85$: update; \textbf{return} \\
\textbf{4.} Mark $p$.\texttt{abstract\_source} = \texttt{missing\_or\_unverified} \\
\hline
\end{tabular}
\caption{Abstract enrichment cascade applied to all 886 full-text queue papers.}
\label{alg:enrichment}
\end{figure}

The cascade raised abstract coverage from 16.6\% (147/886) to 72.6\%
(643/886) before Round~2 screening. Of the 643 enriched abstracts,
590 were sourced from Semantic Scholar, 31 from OpenAlex, 13 from CORE,
and 9 from Crossref or the title-based fallback.
Table~\ref{tab:app-enrichment-stats} summarises the enrichment outcome
by priority band.

\begin{table}[htbp]
\centering
\caption{Abstract enrichment results by priority band.}
\label{tab:app-enrichment-stats}
\small
\setlength{\tabcolsep}{4pt}
\begin{tabular}{c r r r r}
\toprule
\textbf{Band} & \textbf{Papers} & \textbf{With abstract (pre)} &
  \textbf{Gained} & \textbf{Coverage (post)} \\
\midrule
A &  95 & 91 (95.8\%) &   4 &  95 (100.0\%) \\
B &  20 & 18 (90.0\%) &   2 &  20 (100.0\%) \\
C &  39 & 39 (100\%)  &   0 &  39 (100.0\%) \\
D & 152 &  0  (0.0\%) & 118 & 118  (77.6\%) \\
E & 584 &  0  (0.0\%) & 371 & 371  (63.5\%) \\
\midrule
\textbf{Total} & \textbf{886} & \textbf{148 (16.7\%)} &
  \textbf{495} & \textbf{643 (72.6\%)} \\
\bottomrule
\end{tabular}
\end{table}

\section{Quality Assessment Checklist}
\label{app:slr-qa}

The eight-item quality checklist (Table~\ref{tab:app-qa-full}) is applied
to each included primary study during the manual reading phase. Items
QA1--QA4 assess methodological rigour; QA5--QA7 assess empirical soundness
and reproducibility; QA8 specifically targets process model quality reporting,
which is central to PATHCAST's evaluation strategy. Studies scoring below
4/8 are excluded from the synthesis.

\begin{table}[htbp]
\centering
\caption{Quality assessment checklist with scoring guidance.}
\label{tab:app-qa-full}
\small
\setlength{\tabcolsep}{4pt}
\begin{tabular}{c p{5.5cm} p{5.5cm}}
\toprule
\textbf{ID} & \textbf{Criterion} & \textbf{Scoring guidance} \\
\midrule
QA1 & Are the research objectives clearly stated? &
  Score 1 if RQs or hypotheses are explicitly enumerated; 0 otherwise. \\
QA2 & Is the software engineering context described in sufficient detail? &
  Score 1 if the SDLC phase, artefact type, and process scope are defined. \\
QA3 & Is the data source (event log or repository) described reproducibly? &
  Score 1 if dataset origin, collection method, and key statistics are given. \\
QA4 & Is the PM or stochastic technique formally defined? &
  Score 1 if algorithm, parameters, and configuration are specified. \\
QA5 & Are results validated empirically? &
  Score 1 if validation uses real or semi-real data; 0 for simulation-only. \\
QA6 & Are threats to validity discussed? &
  Score 1 if at least internal and external validity are addressed. \\
QA7 & Is the study reproducible (data and/or code available)? &
  Score 1 if a replication package or open dataset is linked or described. \\
QA8 & Are process model quality metrics reported (fitness, precision)? &
  Score 1 if at least one formal quality metric is quantified. \\
\midrule
\multicolumn{2}{l}{\textbf{Total score range}} & 0--8 (threshold: $\geq$4) \\
\bottomrule
\end{tabular}
\end{table}

Current QA status: 76 of 169 studies assessed (full Round~1 set); mean
score = 5.18 (SD = 1.13), median = 5.0 (IQR: 4--6); 5 studies scored
below the threshold and were excluded from the synthesis. The remaining
93 studies (Rounds~2+) are pending manual QA assessment.

\section{Control Paper Set}
\label{app:slr-control}

The 12-paper control set was used to validate the search strings before
execution. Table~\ref{tab:app-control} lists all control papers with their
IC profile and the source through which each was captured.

\begin{table}[htbp]
\centering
\caption{Complete control paper set for search string validation.}
\label{tab:app-control}
\small
\setlength{\tabcolsep}{3pt}
\begin{tabular}{c p{4.0cm} p{3.8cm} l l}
\toprule
\textbf{\#} & \textbf{Paper} & \textbf{Justification} & \textbf{ICs} &
  \textbf{Captured via} \\
\midrule
V1  & Cook \& Wolf (1998)~\cite{cook1998}              & Seminal PM + SE             & IC1, IC2 & ACM \\
V2  & Rubin et al.\ (2007)~\cite{rubin2007}            & First named PM framework for SW & IC1 & Scopus \\
V3  & Poncin et al.\ (2011)~\cite{poncin2011}          & PM + MSR intersection       & IC1, IC4 & IEEE \\
V4  & Caldeira \& Cardoso (2019)                        & PM + SW process assessment  & IC1      & IEEE \\
V5  & Lemos et al.\ (2011)                              & PM for SW management        & IC1      & IEEE \\
V6  & Whittaker \& Thomason (1994)~\cite{whittaker1994} & Markov + SW testing (seminal) & IC2    & Scopus \\
V7  & Tax et al.\ (2017)~\cite{tax2017}                & PPM seminal (BPM baseline)  & IC3      & Control \\
V8  & Rogge-Solti \& Weske (2015)                       & Stochastic + PM forecasting & IC2, IC3 & Scopus \\
V9  & Rozinat et al.\ (2009)                            & PM $\rightarrow$ Simulation  & IC1, IC2 & Control \\
V10 & Gupta (2017)~\cite{gupta2017}                    & PM multi-repository SE      & IC1, IC3 & ACM \\
V11 & Jalote \& Kamma (2021)~\cite{jalote2021}         & IDE logs + Markov chains    & IC1, IC2 & Scopus \\
V12 & Ortu et al.\ (2023)~\cite{ortu2023}              & GitHub Markov + SNA         & IC2, IC4 & Scopus \\
\midrule
\multicolumn{4}{l}{Capture rate from primary queries alone} & 10/12 (83.3\%) \\
\multicolumn{4}{l}{Capture rate including snowballing}      & 12/12 (100\%) \\
\bottomrule
\end{tabular}
\end{table}

Papers V4, V5, V8, and V9 were not in the initial working set but were
recovered through snowballing during the FT phase; they are included in
the final corpus.

\section{SPMF Integration-Level Taxonomy}
\label{app:slr-spmf}

The \textit{Software Process Mining and Forecasting} (SPMF) taxonomy, used
throughout Chapter~\ref{ch:slr}, classifies studies along three orthogonal
dimensions (Table~\ref{tab:app-spmf}).

\begin{table}[htbp]
\centering
\caption{SPMF taxonomy: three dimensions with level definitions.}
\label{tab:app-spmf}
\small
\setlength{\tabcolsep}{4pt}
\begin{tabular}{l l p{7.8cm}}
\toprule
\textbf{Dimension} & \textbf{Level} & \textbf{Definition} \\
\midrule
\multirow{4}{*}{\textbf{Analytical Depth}}
  & AD1 -- Descriptive  & Discovery of as-is process models; conformance checking \\
  & AD2 -- Diagnostic   & Bottleneck detection; rework quantification; root-cause analysis \\
  & AD3 -- Predictive   & Outcome prediction; remaining-time estimation \\
  & AD4 -- Forecasting  & Distributional forecasts; probabilistic completion estimates \\
\midrule
\multirow{5}{*}{\textbf{SDLC Coverage}}
  & SC-R & Requirements and design \\
  & SC-D & Development (coding, commits, pull requests) \\
  & SC-Q & Quality (testing, reviews, defects) \\
  & SC-C & CI/CD and delivery \\
  & SC-O & Operations and maintenance \\
\midrule
\multirow{4}{*}{\textbf{Technique Integration}}
  & TI-L0 & Single technique (PM only, ML only, or stochastic only) \\
  & TI-L1 & Two techniques present in same study without explicit pipeline contract \\
  & TI-L2 & Partial bridge: two techniques linked, but incomplete or missing ML closure \\
  & TI-L3 & Unified pipeline: PM $\rightarrow$ Markov $\rightarrow$ MC $\rightarrow$ ML
             with formal inter-stage contracts \\
\bottomrule
\end{tabular}
\end{table}

The L0--L3 scale in dimension TI maps to the integration-level coding used
in Sections~\ref{sec:slr-rq3} and~\ref{app:slr-top30}. PATHCAST is positioned
at TI-L3, AD4, with SC coverage spanning SC-D, SC-Q, and SC-C.


% =============================================================================
% APPENDIX E — Reading List and Pipeline Artefacts
% =============================================================================
\chapter{Reading List and Pipeline Artefacts}
\label{app:E}

This appendix presents the ranked reading list of 32 primary studies
selected for in-depth synthesis and documents the scripts and persistent
artefacts of the review pipeline.

\section{Multi-Criteria Scoring Rubric}
\label{app:slr-ranking-criteria}

Table~\ref{tab:app-ranking-criteria} presents the 15-point rubric used to
rank studies. Three seminal papers (Cook 1995, Cook 1998, Rubin 2007) are
included regardless of score.

\begin{table}[htbp]
\centering
\caption{Multi-criteria scoring rubric for the top-32 reading list.}
\label{tab:app-ranking-criteria}
\small
\setlength{\tabcolsep}{4pt}
\begin{tabular}{p{7.5cm} c}
\toprule
\textbf{Criterion} & \textbf{Points} \\
\midrule
IC2 present (stochastic modelling in SE)          & +3 \\
IC3 present (forecasting of process metrics)      & +3 \\
IC1 present (process mining in SE)                & +2 \\
IC4 present (repository mining for process)       & +1 \\
PDF locally available                             & +2 \\
Published in 2020 or later                        & +1 \\
High-impact venue (IEEE/ACM/IST/JSS/IS/TSE)      & +1 \\
Empirical study type (case\_study or experiment)  & +1 \\
Contribution type = prediction or simulation      & +1 \\
\midrule
\textbf{Maximum score}                            & \textbf{15} \\
\bottomrule
\end{tabular}
\end{table}

\section{Top-32 Ranked Reading List}
\label{app:slr-top30}

\begin{table}[p]
\centering
\caption{Top-32 reading list with score, IC cluster, and main finding
  summary. Papers marked $\dagger$ are forced-inclusion seminal works.}
\label{tab:app-top30}
\small
\setlength{\tabcolsep}{3pt}
\begin{tabular}{c c p{3.8cm} c c p{5.7cm}}
\toprule
\textbf{Rank} & \textbf{Pts} & \textbf{Authors (year)} &
  \textbf{ICs} & \textbf{PDF} & \textbf{Main finding (summary)} \\
\midrule
 1 & 14 & Joshi \& Kahani (2024)~\cite{joshi2024}
         & IC2,IC3,IC4 & \checkmark
         & RL/MDP for GitHub PR outcome prediction \\
 2 & 14 & Li et al.\ (2020)~\cite{li2020}
         & IC2,IC3 & \checkmark
         & Hybrid SD+DES simulation for industrial project forecasting \\
 3 & 13 & Guinea-Cabrera et al.\ (2025)~\cite{guinea2025}
         & IC1,IC2,IC3 & \checkmark
         & PRIMAD Digital Twin for SDLC (pm4py + Akka + LightGBM), L2 \\
 4 & 13 & Ortu et al.\ (2023)~\cite{ortu2023}
         & IC2,IC4 & \checkmark
         & Fault-insertion/fixing Markov + SNA in Apache projects \\
 5 & 13 & Ben Mesmia et al.\ (2021)~\cite{benmesmia2021}
         & IC2,IC3 & \checkmark
         & Non-Markovian SPN for DevOps workflow duration prediction \\
 6 & 12 & Buliga et al.\ (2025)~\cite{buliga2025}
         & IC1,IC3 & \checkmark
         & Generative AI + business process simulation for what-if \\
 7 & 12 & Bhadra et al.\ (2023)~\cite{bhadra2023}
         & IC2,IC3 & \checkmark
         & SPN model of CI/CD with reliability sensitivity analysis \\
 8 & 12 & Bhadra (2022)~\cite{bhadra2022}
         & IC2,IC3 & \checkmark
         & SPN model of CI/CD pipeline delivery probability \\
 9 & 12 & Massitela et al.\ (2018)~\cite{massitela2018}
         & IC2,IC3 & \checkmark
         & Stochastic Automata Network for Waterfall team estimation \\
10 & 11 & Incerto et al.\ (2025)~\cite{incerto2025}
         & IC1,IC2 & \checkmark
         & VLMC stochastic conformance; outperforms 18 state-of-art techniques \\
11 & 11 & Jo \& Kwon (2024)~\cite{jo2024}
         & IC2,IC4 & \checkmark
         & Probabilistic model checking (CTL) on GitHub repos \\
12 & 11 & Jo et al.\ (2023)~\cite{jo2023}
         & IC2,IC4 & \checkmark
         & DTMC collaboration patterns in GitHub (5 repos) \\
13 & 11 & L\'opez-Pintado \& Dumas (2023)~\cite{lopezpintado2023}
         & IC1,IC2 & ---
         & Probabilistic resource calendars discovered from event logs \\
14 & 11 & Caldeira et al.\ (2022)~\cite{caldeira2022}
         & IC1,IC3 & \checkmark
         & Process complexity $\to$ cyclomatic complexity ($r \approx 0.43$) \\
15 & 11 & Jalote \& Kamma (2021)~\cite{jalote2021}
         & IC1,IC2 & \checkmark
         & Markov chains on IDE logs for programmer productivity \\
16 & 11 & Lunesu et al.\ (2021)~\cite{lunesu2021}
         & IC2,IC3 & \checkmark
         & Monte Carlo risk assessment for agile software development \\
17 & 11 & Nafreen et al.\ (2020)~\cite{nafreen2020}
         & IC2,IC3 & ---
         & SRGM + defect tracking for hybrid reliability forecasting \\
18 & 11 & Gupta (2017)~\cite{gupta2017}
         & IC1,IC3,IC4 & ---
         & Inductive miner + predictive analytics for SW maintenance \\
19 & 10 & Salmani et al.\ (2022)
         & IC1,IC3 & \checkmark
         & Process mining for requirements engineering in MDPs \\
20 & 10 & Jeon et al.\ (2015)~\cite{jeon2015}
         & IC2,IC3,IC4 & ---
         & Markov-based probabilistic risk prediction from repo data \\
21 & 10 & Tian \& Noore (2005)~\cite{tian2005}
         & IC2,IC3 & ---
         & Markov defect-removal modelling under code churn \\
22 & 10 & Cook \& Wolf (1995)$^\dagger$~\cite{cook1995}
         & IC1,IC2 & \checkmark
         & First automated process discovery from event traces \\
23 &  9 & Pourbafrani et al.\ (2021)~\cite{pourbafrani2021}
         & IC1,IC3 & \checkmark
         & SIMPT: process tree + interactive what-if simulation tool \\
24 &  9 & Tyagi et al.\ (2021)~\cite{tyagi2021}
         & IC2,IC3 & ---
         & State-space Markov framework for agile planning and tracking \\
25 &  9 & Razi et al.\ (2019)
         & IC1,IC2 & \checkmark
         & Conformance checking for SDLC process compliance \\
26 &  9 & Krismayer et al.\ (2019)
         & IC1,IC4 & \checkmark
         & Declarative PM for repository-based process monitoring \\
27 &  9 & Baskoro et al.\ (2019)
         & IC3,IC4 & ---
         & ML-based effort prediction from repository data \\
28 &  9 & (Anon.\ 2017)
         & IC2,IC3 & ---
         & Probabilistic defect prediction for software product lines \\
29 &  9 & Magennis (2015)~\cite{magennis2015}
         & IC2,IC3 & ---
         & Monte Carlo cycle-time economic analysis \\
30 &  9 & Washizaki et al.\ (2015)~\cite{washizaki2015}
         & IC2,IC3 & ---
         & GSRM release-time prediction for open-source projects \\
31 &  9 & Cook \& Wolf (1998)$^\dagger$~\cite{cook1998}
         & IC1,IC2 & ---
         & Markov/FSA/NN process discovery from event-based data \\
32 &  2 & Rubin et al.\ (2007)$^\dagger$~\cite{rubin2007}
         & IC1 & \checkmark
         & First named PM framework for SDLC (ProM + SCM logs) \\
\bottomrule
\end{tabular}
\end{table}

\section{Review Pipeline Scripts}
\label{app:slr-pipeline}

The review was conducted using a custom Python pipeline (\texttt{pipeline/})
operating on a structured CSV data store. Table~\ref{tab:app-pipeline}
documents each script with its function and main invocation command.

\begin{table}[htbp]
\centering
\caption{Review pipeline scripts and their roles.}
\label{tab:app-pipeline}
\small
\setlength{\tabcolsep}{3pt}
\begin{tabular}{p{4.2cm} p{4.5cm} p{4.0cm}}
\toprule
\textbf{Script} & \textbf{Function} & \textbf{Main command} \\
\midrule
\texttt{pipeline/dedup.py} &
  Cross-source deduplication (DOI + fuzzy title) &
  \texttt{python main.py dedup} \\[2pt]
\texttt{pipeline/enrich.py} &
  Abstract enrichment cascade (4 sources) &
  \texttt{python main.py enrich} \\[2pt]
\texttt{pipeline/screening.py} &
  T/A LLM screening via Batches API &
  \texttt{python main.py screen -{}-run} \\[2pt]
\texttt{pipeline/fulltext.py} &
  FT LLM screening + manual review &
  \texttt{python main.py fulltext -{}-run} \\[2pt]
\texttt{pipeline/pdf\_downloader.py} &
  PDF retrieval (Unpaywall, S2, OA) &
  \texttt{python main.py pdf} \\[2pt]
\texttt{pipeline/extract\_prep.py} &
  S2 metadata enrichment + PDF copy &
  \texttt{python pipeline/} \texttt{extract\_prep.py} \\[2pt]
\texttt{pipeline/extract\_llm.py} &
  Structured data extraction (LLM, 169 papers) &
  \texttt{python pipeline/} \texttt{extract\_llm.py -{}-run} \\[2pt]
\texttt{pipeline/synth\_llm.py} &
  Synthesis paragraph generation per IC cluster &
  \texttt{python pipeline/} \texttt{synth\_llm.py} \\[2pt]
\texttt{pipeline/finalization.py} &
  PRISMA snapshot, EC5 closure, artefact export &
  \texttt{python main.py finalize} \\
\bottomrule
\end{tabular}
\end{table}

\section{Persistent Artefact Registry}
\label{app:slr-artefacts}

Table~\ref{tab:app-artefacts} lists the persistent artefacts produced by the
review, with their role, location, and current state.

\begin{table}[htbp]
\centering
\caption{Persistent artefacts of the review.}
\label{tab:app-artefacts}
\small
\setlength{\tabcolsep}{3pt}
\begin{tabular}{p{6.5cm} p{3.0cm} p{2.8cm}}
\toprule
\textbf{File} & \textbf{Role} & \textbf{State} \\
\midrule
\texttt{results/screening/} \texttt{ft\_screening\_results.csv}
  & 886-paper decision log & Final \\[3pt]
\texttt{results/extraction/} \texttt{extraction\_template.csv}
  & 169-paper extraction data & 169/169 \\[3pt]
\texttt{results/final\_review/} \texttt{included\_studies\_current.csv}
  & Included study set & 169 rows \\[3pt]
\texttt{results/final\_review/} \texttt{top30\_reading\_list.csv}
  & Ranked reading list & 32 rows \\[3pt]
\texttt{results/final\_review/} \texttt{cap3\_synthesis\_v2.tex}
  & LLM synthesis paragraphs & For review \\[3pt]
\texttt{results/final\_review/} \texttt{missing\_references.bib}
  & BibTeX for 32 missing keys & Ready to merge \\[3pt]
\texttt{results/final\_review/} \texttt{pending\_inaccessible\_closed.csv}
  & EC5 closure record & 148 rows \\[3pt]
\texttt{results/final\_review/} \texttt{prisma\_summary\_current.txt}
  & PRISMA snapshot & Final \\[3pt]
\texttt{results/} \texttt{SLR\_EXECUCAO\_ATE\_AGORA.md}
  & Execution log (all sessions) & Up to date \\
\bottomrule
\end{tabular}
\end{table}

% The \appendix command causes subsequent \chapter{} calls to be labelled
% as APPENDIX A, B, C, etc. per ABNT NBR 14724.
% =============================================================================


% =============================================================================
% APPENDIX A — Search Protocol: Search Strings
% =============================================================================
\chapter{Search Protocol: Search Strings}
\label{app:A}

This appendix presents the four search strings executed across the five
digital libraries during the search phase of the PATHCAST SLR. Each string
was adapted to the syntax and field conventions of the target database; the
block structure (P and I) is preserved across all adaptations. Strings S1--S3
correspond to the principal, complementary (MSR), and frontier (stochastic
$\wedge$ PM) blocks, respectively; S4 is a supplementary string executed in
Scopus only, targeting Markov chain models in software testing.

\section{Search String Summary by Database}
\label{app:slr-strings}

\begin{table}[htbp]
\centering
\caption{Search strings by database and field scope.}
\label{tab:app-strings}
\small
\setlength{\tabcolsep}{4pt}
\begin{tabular}{p{2.4cm} c p{9.0cm}}
\toprule
\textbf{Database} & \textbf{Strings} & \textbf{Field / syntax notes} \\
\midrule
Scopus       & S1, S2, S3, S4 & \texttt{TITLE-ABS-KEY(...)};
                                  \texttt{DOCTYPE(ar OR cp)};
                                  \texttt{LANGUAGE(english)};
                                  1994--2026 \\
IEEE Xplore  & S1 (split)     & All Metadata; Boolean AND of sub-blocks;
                                  Conference + Journals; 1994--2026 \\
Web of Science & S1, S2, S3   & Topic \texttt{TS=(...)}; Article, Proceedings;
                                  English; 1994--2026 \\
SpringerLink & S1 (simplified)& Title + Abstract; 8 separate exports;
                                  Chapter + Article; English \\
ACM DL       & S1, S2         & Abstract; 2 queries; date 1994--2026 \\
\bottomrule
\end{tabular}
\end{table}

\section{String S1 — Principal ($\text{P} \wedge \text{I}$)}
\label{app:slr-s1}

Block P (software engineering context):

\begin{quote}
\small\ttfamily
("software engineering" OR "software development" OR "software process"
OR "software lifecycle" OR "software life cycle" OR "SDLC" OR "development
lifecycle" OR "development life cycle" OR "DevOps" OR "CI/CD" OR
"continuous integration" OR "continuous delivery" OR "agile development"
OR "agile process" OR "issue tracking" OR "bug tracking" OR "version
control" OR "code review" OR "pull request" OR "software project"
OR "software repository" OR "software maintenance" OR "software
evolution")
\end{quote}

Block I (techniques):

\begin{quote}
\small\ttfamily
("process mining" OR "process discovery" OR "conformance checking" OR
"workflow mining" OR "predictive process monitoring" OR "process prediction"
OR "remaining time prediction" OR "process forecasting" OR "event log" OR
"process simulation" OR "Monte Carlo simulation" OR "Markov chain" OR
"Markov model" OR "stochastic process model" OR "transition probability"
OR "absorbing chain" OR "lead time prediction" OR "cycle time prediction")
\end{quote}

\section{String S2 — Complementary (MSR $\rightarrow$ Process Modeling)}
\label{app:slr-s2}

\begin{quote}
\small\ttfamily
("mining software repositories" OR "software repository mining" OR
"GitHub mining" OR "Jira mining") AND ("process model" OR "Markov" OR
"Monte Carlo" OR "stochastic" OR "lead time" OR "cycle time" OR "throughput"
OR "transition matrix")
\end{quote}

\section{String S3 — Frontier (Stochastic $\wedge$ PM $\wedge$ SE)}
\label{app:slr-s3}

\begin{quote}
\small\ttfamily
("Monte Carlo" OR "Markov chain" OR "stochastic simulation" OR
"probabilistic forecast" OR "transition matrix") AND ("process mining" OR
"event log" OR "workflow mining" OR "process model" OR "business process")
AND ("software development" OR "software engineering" OR "software process"
OR "DevOps" OR "SDLC" OR "issue tracker" OR "commit log" OR "build pipeline")
\end{quote}

\section{String S4 — Markov Chain Software Testing (Scopus supplement)}
\label{app:slr-s4}

\begin{quote}
\small\ttfamily
("Markov chain" OR "Markov model" OR "hidden Markov" OR "stochastic automata")
AND ("software testing" OR "test coverage" OR "software reliability" OR
"defect prediction" OR "fault prediction")
\end{quote}


% =============================================================================
% APPENDIX B — LLM Screening Prompts
% =============================================================================
\chapter{LLM Screening Prompts}
\label{app:B}

The review employed \textit{claude-haiku-4-5-20251001} as primary screener in
two phases: title/abstract (T/A) and full-text (FT). Each phase uses a
\emph{system prompt} that encodes the scope, criteria, and output schema,
paired with a per-paper \emph{user prompt} that supplies the bibliographic
data. All prompts are version-controlled in
\texttt{config/screening\_criteria.py} in the review repository.

\section{Phase 1: Title/Abstract Screening}
\label{app:slr-ta-prompt}

\subsection*{System Prompt (T/A Phase)}

\begin{quote}
\small\ttfamily
You are an expert in systematic literature review (SLR) in the area of
process mining and software engineering. Your task is to perform the
title/abstract (T/A) screening for SLR PATHCAST.

\textbf{SLR PATHCAST scope:}
Intersection of (1) process mining, (2) stochastic modelling,
(3) forecasting/prediction, and (4) software development processes.
Focus: studies that analyse, discover, or predict software process
characteristics using event logs, repositories, or stochastic models.

\textbf{INCLUSION CRITERIA} (at least ONE must be satisfied to include):
IC1 -- Process Mining in Software: applies PM (discovery, conformance,
workflow mining, event log analysis, PPM) to SW artefacts (commits,
issues, PRs, CI/CD).
IC2 -- Stochastic Modelling of SW Processes: uses Markov chains, Monte
Carlo, stochastic Petri nets, or transition matrices to model SW dev
processes.
IC3 -- Process Metric Forecasting: predicts lead time, cycle time,
remaining time, throughput, or defect rates in SW processes using
data or event logs.
IC4 -- Repository Mining for Process: mines repositories (GitHub, Jira,
VCS) to discover or improve SW development process models.

\textbf{EXCLUSION CRITERIA} (any ONE is sufficient to exclude):
EC1 -- Non-SW domain: applied to healthcare, manufacturing, finance,
etc., with no link to SW processes.
EC2 -- SW as tool: "software" is the implementation platform, not the
process under study.
EC3 -- Algorithm without SW application: purely theoretical method,
no evaluation in SW processes.
EC4 -- Out-of-scope secondary study: survey/SLR not specifically
addressing PM or stochastic methods in SE.

\textbf{UNCERTAIN CASES (maybe):}
Relevant title without abstract; generic BPM article with possible SW
application; ambiguous terminology.

\textbf{STRUCTURED FIELDS} (use only the allowed enums):
evidence\_tags, software\_context, stochastic\_method,
forecast\_target, process\_data\_source, confidence.

Respond ONLY with valid JSON, no additional text.
\end{quote}

\subsection*{User Prompt Template (T/A Phase)}

\begin{quote}
\small\ttfamily
Evaluate the paper below for SLR PATHCAST.

\textbf{Title:} \{title\}
\textbf{Abstract:} \{abstract\}
\textbf{Venue/Type:} \{venue\} | \{doc\_type\} | \{year\}
\textbf{Source database:} \{source\_db\}
\textbf{Abstract origin:} \{abstract\_source\}

Respond with JSON in the exact format:
\{
  "decision": "include" | "exclude" | "maybe",
  "rationale": "<1-2 sentences justifying the decision>",
  "matched\_ic": ["IC1", "IC2"],
  "matched\_ec": ["EC1"],
  "evidence\_tags": ["process\_mining", "forecasting"],
  "software\_context": "software\_development\_process",
  "stochastic\_method": "markov\_chain",
  "forecast\_target": "remaining\_time",
  "process\_data\_source": "event\_logs",
  "confidence": "medium"
\}

Rules:
- "include" if $\geq$1 IC met and no decisive EC
- "exclude" if an EC clearly applies, or paper clearly meets no IC
- "maybe" if genuine doubt (title without abstract, domain ambiguity)
- matched\_ic and matched\_ec must be [] when not applicable
\end{quote}

\section{Phase 2: Full-Text Screening}
\label{app:slr-ft-prompt}

\subsection*{System Prompt (FT Phase)}

The FT system prompt is identical to the T/A system prompt above, with
three additional instructions that encode the stricter standard of the
full-text phase:

\begin{quote}
\small\ttfamily
THIS IS THE FULL-TEXT PHASE -- apply criteria rigorously:
- If the abstract clearly confirms $\geq$1 IC with no decisive EC
  $\rightarrow$ "include"
- If the abstract clearly confirms no IC applies, or an EC applies
  $\rightarrow$ "exclude"
- Use "pending" ONLY if the abstract is genuinely insufficient for
  a decision (absent or extremely short) AND the title suggests
  marginal relevance

When in doubt between include and exclude, prefer exclude -- this
phase requires confidence.

YOU MUST NEVER invent, complete, or infer a missing abstract.
\end{quote}

\subsection*{User Prompt Template (FT Phase)}

The FT user prompt extends the T/A template with a \textbf{T/A context block}
that supplies the previous-phase decision, rationale, and matched ICs:

\begin{quote}
\small\ttfamily
Evaluate the paper below for SLR PATHCAST -- FULL-TEXT screening (phase 2).

\textbf{Title:} \{title\}
\textbf{Abstract:} \{abstract\}
\textbf{Venue/Type:} \{venue\} | \{doc\_type\} | \{year\}
\textbf{Source database:} \{source\_db\}
\textbf{Abstract origin:} \{abstract\_source\}

\textbf{T/A phase context:}
  - Previous decision: \{ta\_decision\}
  - Previous rationale: \{ta\_rationale\}
  - ICs identified in phase 1: \{ta\_matched\_ic\}

Re-evaluate rigorously for the full-text phase.
Respond with JSON in the exact format:
\{
  "decision": "include" | "exclude" | "pending",
  [... same fields as T/A phase ...]
\}
\end{quote}

\section{Phase 2b: Screen-Blanks Pass}
\label{app:slr-blanks-prompt}

A third pass targeted the 186 papers without a decision after both prior
phases (primarily Band~E papers with no available abstract). The screen-blanks
pass reuses the same FT system and user prompts, with the abstract field
populated by the placeholder ``Abstract not available'' and the T/A context
block showing \texttt{ta\_decision~=~maybe} and
\texttt{ta\_matched\_ic~=~none}. Papers for which the LLM returned
\texttt{pending} and whose title analysis suggested low relevance were
subsequently closed under EC5 (full text inaccessible; $n = 148$).

\section{Model Parameters per Screening Phase}
\label{app:slr-prompt-params}

Table~\ref{tab:app-prompt-params} summarises the model parameters used
across all three screening phases and the data extraction step.

\begin{table}[htbp]
\centering
\caption{LLM parameters used at each screening phase.}
\label{tab:app-prompt-params}
\small
\setlength{\tabcolsep}{4pt}
\begin{tabular}{l l r r l}
\toprule
\textbf{Phase} & \textbf{Model} & \textbf{Max tok.} &
  \textbf{Batch} & \textbf{API} \\
\midrule
T/A screening      & claude-haiku-4-5-20251001 & 512  & 500 & Anthropic Batches \\
FT screening       & claude-haiku-4-5-20251001 & 512  & 100 & Anthropic Batches \\
Screen-blanks pass & claude-haiku-4-5-20251001 & 512  & 186 & Anthropic Batches \\
Data extraction    & claude-haiku-4-5-20251001 & 1024 & 100 & Anthropic Batches \\
Synthesis (ch.~3)  & claude-sonnet-4-6         & 2048 & --- & Anthropic Messages \\
\bottomrule
\end{tabular}
\end{table}


% =============================================================================
% APPENDIX C — Data Extraction Form
% =============================================================================
\chapter{Data Extraction Form}
\label{app:C}

Data extraction for all 169 included studies was performed with
\textit{claude-haiku-4-5-20251001} via the Anthropic Batches API. Papers
with a locally available PDF had their text pre-extracted with
\texttt{pdfplumber} (first 8 pages, maximum 6,000 characters). Papers
without a PDF used the enriched abstract from Semantic Scholar.
The prompt below was submitted once per paper.

\section{Extraction Prompt}
\label{app:slr-extraction-prompt}

\begin{quote}
\small\ttfamily
You are a researcher performing data extraction for SLR PATHCAST on
Process Mining and Stochastic Modelling in Software Processes.

\textbf{Title:} \{title\}
\textbf{Authors:} \{authors\}
\textbf{Year:} \{year\} | \textbf{Venue:} \{venue\}
\textbf{Inclusion criteria met:} \{ics\}
\{content\_block\}

Return ONLY a valid JSON object with exactly these fields:
\{
  "research\_question": "main research question (1-2 sentences)",
  "study\_type": "case\_study | experiment | survey | tool |
                 theoretical | simulation | mixed",
  "research\_contribution": "discovery | conformance | enhancement |
       prediction | simulation | framework | hybrid",
  "pm\_technique": "alpha | inductive\_miner | heuristic |
       conformance\_checking | social\_network | declarative | other | none",
  "stochastic\_technique": "markov\_chain | stochastic\_petri\_net |
       monte\_carlo | system\_dynamics | other | none",
  "software\_artifact": "commits | issues | ci\_cd | ide\_logs |
       vcs | jira | github | other | none",
  "software\_process": "development | testing | maintenance |
       code\_review | bug\_fixing | deployment | requirements | other",
  "dataset\_source": "open\_source | industrial | academic |
                     synthetic | not\_specified",
  "dataset\_public": "yes | no | partial | not\_mentioned",
  "tool\_used": "tool names (e.g.\ ProM, pm4py, Disco) or not\_mentioned",
  "main\_finding": "main result/contribution (2-3 sentences)",
  "limitations": "reported limitations (1-2 sentences) or not\_mentioned",
  "replication\_package": "yes | no | partial | not\_mentioned"
\}

Respond ONLY with JSON, no additional text.
\end{quote}

\section{Extraction Form Fields}
\label{app:slr-extraction-fields}

Table~\ref{tab:app-extraction-fields} describes all fields of the structured
extraction form stored in \texttt{extraction\_template.csv}.

\begin{table}[htbp]
\centering
\caption{Data extraction form: all fields in \texttt{extraction\_template.csv}.}
\label{tab:app-extraction-fields}
\footnotesize
\setlength{\tabcolsep}{3pt}
\begin{tabular}{p{3.2cm} p{4.4cm} p{4.6cm}}
\toprule
\textbf{Field} & \textbf{Description} & \textbf{Allowed values / format} \\
\midrule
\multicolumn{3}{l}{\textit{Identification}} \\
\texttt{internal\_id}   & 8-char hash (dedup key)      & --- \\
\texttt{title}          & Paper title                  & free text \\
\texttt{doi}            & Digital Object Identifier    & URI \\
\texttt{year}           & Publication year             & YYYY \\
\texttt{source\_db}     & Source database              & scopus, ieee, acm, springer, wos, snowball \\
\texttt{ft\_matched\_ic}& Inclusion criteria met       & IC1$|$IC2$|$IC3$|$IC4 \\
\texttt{ft\_screened\_by}& Who made FT decision        & llm, manual, llm\_pdf \\
\midrule
\multicolumn{3}{l}{\textit{Enriched metadata (Semantic Scholar)}} \\
\texttt{authors}        & Author list                  & semicolon-separated \\
\texttt{venue}          & Publication venue            & free text \\
\texttt{journal\_name}  & Journal/conference name      & free text \\
\texttt{abstract}       & Paper abstract               & free text \\
\texttt{pdf\_available} & PDF locally present          & yes / no \\
\midrule
\multicolumn{3}{l}{\textit{LLM-extracted fields}} \\
\texttt{research\_question} & Main RQ               & 1--2 sentences \\
\texttt{study\_type}        & Study design          & case\_study $|$ experiment $|$ survey $|$ tool $|$ theoretical $|$ simulation $|$ mixed \\
\texttt{research\_contribution} & Contribution type & discovery $|$ conformance $|$ enhancement $|$ prediction $|$ simulation $|$ framework $|$ hybrid \\
\texttt{pm\_technique}      & PM algorithm used     & alpha $|$ inductive\_miner $|$ heuristic $|$ conformance\_checking $|$ social\_network $|$ declarative $|$ other $|$ none \\
\texttt{stochastic\_technique} & Stochastic method  & markov\_chain $|$ stochastic\_petri\_net $|$ monte\_carlo $|$ system\_dynamics $|$ other $|$ none \\
\texttt{software\_artifact} & Data source           & commits $|$ issues $|$ ci\_cd $|$ ide\_logs $|$ vcs $|$ jira $|$ github $|$ other $|$ none \\
\texttt{software\_process}  & SDLC phase            & development $|$ testing $|$ maintenance $|$ code\_review $|$ bug\_fixing $|$ deployment $|$ requirements $|$ other \\
\texttt{dataset\_source}    & Dataset provenance    & open\_source $|$ industrial $|$ academic $|$ synthetic $|$ not\_specified \\
\texttt{dataset\_public}    & Data availability     & yes $|$ no $|$ partial $|$ not\_mentioned \\
\texttt{tool\_used}         & Tools/platforms       & free text \\
\texttt{main\_finding}      & Key result            & 2--3 sentences \\
\texttt{limitations}        & Study limitations     & 1--2 sentences \\
\texttt{replication\_package} & Artefacts available & yes $|$ no $|$ partial $|$ not\_mentioned \\
\midrule
\multicolumn{3}{l}{\textit{Manual fields (reading phase)}} \\
\texttt{quality\_score}     & QA checklist score (0--8) & numeric \\
\texttt{extraction\_notes}  & Free-form reviewer note   & free text \\
\bottomrule
\end{tabular}
\end{table}

\section{Structured Output Controlled Vocabulary}
\label{app:slr-schema}

Table~\ref{tab:app-enum-schema} defines the complete controlled vocabulary
for the structured fields returned by the LLM at both screening phases.
These enums are version-controlled in \texttt{config/screening\_criteria.py}
and were not changed between screening rounds to ensure consistency.

\begin{table}[htbp]
\centering
\caption{Controlled vocabulary (enums) for LLM structured output.}
\label{tab:app-enum-schema}
\small
\setlength{\tabcolsep}{4pt}
\begin{tabular}{p{3.2cm} p{9.4cm}}
\toprule
\textbf{Field} & \textbf{Allowed values} \\
\midrule
\texttt{decision} (T/A) &
  \texttt{include}, \texttt{exclude}, \texttt{maybe} \\
\texttt{decision} (FT) &
  \texttt{include}, \texttt{exclude}, \texttt{pending} \\
\texttt{software\_context} &
  \texttt{software\_development\_process}, \texttt{repository\_mining},
  \texttt{ci\_cd}, \texttt{issue\_bug\_workflow}, \texttt{software\_testing},
  \texttt{requirements\_engineering}, \texttt{software\_project\_management},
  \texttt{unclear}, \texttt{not\_software\_process} \\
\texttt{stochastic\_method} &
  \texttt{none}, \texttt{markov\_chain}, \texttt{hidden\_markov\_model},
  \texttt{monte\_carlo}, \texttt{stochastic\_petri\_net},
  \texttt{bayesian\_model}, \texttt{probabilistic\_model},
  \texttt{simulation}, \texttt{queueing\_model},
  \texttt{other\_stochastic}, \texttt{unclear} \\
\texttt{forecast\_target} &
  \texttt{none}, \texttt{lead\_time}, \texttt{cycle\_time},
  \texttt{remaining\_time}, \texttt{throughput}, \texttt{defect\_rate},
  \texttt{build\_outcome}, \texttt{reliability},
  \texttt{completion\_time}, \texttt{other\_process\_metric},
  \texttt{unclear} \\
\texttt{process\_data\_source} &
  \texttt{none}, \texttt{event\_logs}, \texttt{version\_control},
  \texttt{issue\_tracker}, \texttt{pull\_requests}, \texttt{ci\_cd\_logs},
  \texttt{software\_repository\_mixed}, \texttt{synthetic\_data},
  \texttt{simulated\_process}, \texttt{survey\_or\_secondary},
  \texttt{unclear} \\
\texttt{confidence} &
  \texttt{low}, \texttt{medium}, \texttt{high} \\
\texttt{evidence\_tags} &
  \texttt{process\_mining}, \texttt{software\_process},
  \texttt{repository\_mining}, \texttt{stochastic\_modeling},
  \texttt{forecasting}, \texttt{event\_log}, \texttt{version\_control},
  \texttt{issue\_tracking}, \texttt{pull\_requests}, \texttt{ci\_cd},
  \texttt{markov}, \texttt{hidden\_markov\_model}, \texttt{monte\_carlo},
  \texttt{stochastic\_petri\_net}, \texttt{bayesian\_model},
  \texttt{simulation}, \texttt{lead\_time}, \texttt{cycle\_time},
  \texttt{remaining\_time}, \texttt{throughput}, \texttt{defect\_prediction},
  \texttt{build\_prediction}, \texttt{reliability},
  \texttt{insufficient\_abstract}, \texttt{full\_text\_required} \\
\bottomrule
\end{tabular}
\end{table}


% =============================================================================
% APPENDIX D — Assessment and Classification Instruments
% =============================================================================
\chapter{Assessment and Classification Instruments}
\label{app:D}

This appendix gathers the instruments used to classify, enrich, and evaluate
papers throughout the screening and selection process: the priority band
taxonomy, the abstract enrichment cascade, the quality assessment checklist,
the control paper set, and the SPMF integration-level taxonomy.

\section{Priority Bands for the Full-Text Review Queue}
\label{app:slr-bands}

Prior to full-text screening, the 886-paper queue was stratified into five
priority bands (A--E) to optimise resource allocation. Abstract enrichment
and LLM screening were applied in band order.
Table~\ref{tab:app-bands} defines the classification rules.

\begin{table}[htbp]
\centering
\caption{Priority band definitions for the full-text review queue.}
\label{tab:app-bands}
\small
\setlength{\tabcolsep}{4pt}
\begin{tabular}{c r p{8.5cm}}
\toprule
\textbf{Band} & \textbf{$n$} & \textbf{Classification rule} \\
\midrule
A &  95 & T/A decision = \texttt{include} and $\geq$2 distinct ICs matched \\
B &  20 & T/A decision = \texttt{include} and exactly 1 IC matched \\
C &  39 & T/A decision = \texttt{maybe} and abstract available \\
D & 152 & T/A decision = \texttt{maybe}, at least 1 IC signal in title, no abstract \\
E & 584 & T/A decision = \texttt{maybe}, no IC signal in title, no abstract \\
\midrule
\textbf{Total} & \textbf{886} & \\
\bottomrule
\end{tabular}
\end{table}

Band~D papers with a DOI were prioritised for manual PDF retrieval. Band~E
papers without a DOI and without an abstract were closed as EC5 (full text
inaccessible) at the end of the screen-blanks pass.

\section{Abstract Enrichment Cascade}
\label{app:slr-enrichment}

To increase the number of papers assessable by the LLM, a four-source
cascade was executed prior to full-text screening.
Algorithm~\ref{alg:enrichment} describes the procedure.

\begin{figure}[htbp]
\centering
\small
\begin{tabular}{l}
\hline
\textbf{Algorithm 1} \quad Abstract Enrichment Cascade \\
\hline
\textbf{Input:} Paper record $p$ with \texttt{doi} and/or \texttt{title} \\
\textbf{Output:} Updated $p$ with \texttt{abstract}, \texttt{authors}, \texttt{venue} \\
\hline
\textbf{1.} If $p$.\texttt{abstract} is non-empty $\rightarrow$ \textbf{return} $p$ \\
\textbf{2.} If $p$.\texttt{doi} is present: \\
\quad\textbf{2a.} Query \textbf{Semantic Scholar} Batch API (\texttt{DOI:\{doi\}}) \\
\quad\quad If match found: update \texttt{abstract}, \texttt{authors}, \texttt{venue}; \textbf{return} \\
\quad\textbf{2b.} Query \textbf{OpenAlex} (\texttt{filter=doi:\{doi\}}) \\
\quad\quad If match found: update \texttt{abstract}; \textbf{return} \\
\quad\textbf{2c.} Query \textbf{CORE} (\texttt{/search/works?q=doi:\{doi\}}) \\
\quad\quad If match found: update \texttt{abstract}; \textbf{return} \\
\quad\textbf{2d.} Query \textbf{Crossref} (\texttt{/works/\{doi\}}) \\
\quad\quad If abstract present: update; \textbf{return} \\
\textbf{3.} If $p$.\texttt{doi} is absent: \\
\quad\textbf{3a.} Query \textbf{Semantic Scholar} by title (fuzzy) \\
\quad\quad If similarity $\geq 0.85$: update; \textbf{return} \\
\textbf{4.} Mark $p$.\texttt{abstract\_source} = \texttt{missing\_or\_unverified} \\
\hline
\end{tabular}
\caption{Abstract enrichment cascade applied to all 886 full-text queue papers.}
\label{alg:enrichment}
\end{figure}

The cascade raised abstract coverage from 16.6\% (147/886) to 72.6\%
(643/886) before Round~2 screening. Of the 643 enriched abstracts,
590 were sourced from Semantic Scholar, 31 from OpenAlex, 13 from CORE,
and 9 from Crossref or the title-based fallback.
Table~\ref{tab:app-enrichment-stats} summarises the enrichment outcome
by priority band.

\begin{table}[htbp]
\centering
\caption{Abstract enrichment results by priority band.}
\label{tab:app-enrichment-stats}
\small
\setlength{\tabcolsep}{4pt}
\begin{tabular}{c r r r r}
\toprule
\textbf{Band} & \textbf{Papers} & \textbf{With abstract (pre)} &
  \textbf{Gained} & \textbf{Coverage (post)} \\
\midrule
A &  95 & 91 (95.8\%) &   4 &  95 (100.0\%) \\
B &  20 & 18 (90.0\%) &   2 &  20 (100.0\%) \\
C &  39 & 39 (100\%)  &   0 &  39 (100.0\%) \\
D & 152 &  0  (0.0\%) & 118 & 118  (77.6\%) \\
E & 584 &  0  (0.0\%) & 371 & 371  (63.5\%) \\
\midrule
\textbf{Total} & \textbf{886} & \textbf{148 (16.7\%)} &
  \textbf{495} & \textbf{643 (72.6\%)} \\
\bottomrule
\end{tabular}
\end{table}

\section{Quality Assessment Checklist}
\label{app:slr-qa}

The eight-item quality checklist (Table~\ref{tab:app-qa-full}) is applied
to each included primary study during the manual reading phase. Items
QA1--QA4 assess methodological rigour; QA5--QA7 assess empirical soundness
and reproducibility; QA8 specifically targets process model quality reporting,
which is central to PATHCAST's evaluation strategy. Studies scoring below
4/8 are excluded from the synthesis.

\begin{table}[htbp]
\centering
\caption{Quality assessment checklist with scoring guidance.}
\label{tab:app-qa-full}
\small
\setlength{\tabcolsep}{4pt}
\begin{tabular}{c p{5.5cm} p{5.5cm}}
\toprule
\textbf{ID} & \textbf{Criterion} & \textbf{Scoring guidance} \\
\midrule
QA1 & Are the research objectives clearly stated? &
  Score 1 if RQs or hypotheses are explicitly enumerated; 0 otherwise. \\
QA2 & Is the software engineering context described in sufficient detail? &
  Score 1 if the SDLC phase, artefact type, and process scope are defined. \\
QA3 & Is the data source (event log or repository) described reproducibly? &
  Score 1 if dataset origin, collection method, and key statistics are given. \\
QA4 & Is the PM or stochastic technique formally defined? &
  Score 1 if algorithm, parameters, and configuration are specified. \\
QA5 & Are results validated empirically? &
  Score 1 if validation uses real or semi-real data; 0 for simulation-only. \\
QA6 & Are threats to validity discussed? &
  Score 1 if at least internal and external validity are addressed. \\
QA7 & Is the study reproducible (data and/or code available)? &
  Score 1 if a replication package or open dataset is linked or described. \\
QA8 & Are process model quality metrics reported (fitness, precision)? &
  Score 1 if at least one formal quality metric is quantified. \\
\midrule
\multicolumn{2}{l}{\textbf{Total score range}} & 0--8 (threshold: $\geq$4) \\
\bottomrule
\end{tabular}
\end{table}

Current QA status: 76 of 169 studies assessed (full Round~1 set); mean
score = 5.18 (SD = 1.13), median = 5.0 (IQR: 4--6); 5 studies scored
below the threshold and were excluded from the synthesis. The remaining
93 studies (Rounds~2+) are pending manual QA assessment.

\section{Control Paper Set}
\label{app:slr-control}

The 12-paper control set was used to validate the search strings before
execution. Table~\ref{tab:app-control} lists all control papers with their
IC profile and the source through which each was captured.

\begin{table}[htbp]
\centering
\caption{Complete control paper set for search string validation.}
\label{tab:app-control}
\small
\setlength{\tabcolsep}{3pt}
\begin{tabular}{c p{4.0cm} p{3.8cm} l l}
\toprule
\textbf{\#} & \textbf{Paper} & \textbf{Justification} & \textbf{ICs} &
  \textbf{Captured via} \\
\midrule
V1  & Cook \& Wolf (1998)~\cite{cook1998}              & Seminal PM + SE             & IC1, IC2 & ACM \\
V2  & Rubin et al.\ (2007)~\cite{rubin2007}            & First named PM framework for SW & IC1 & Scopus \\
V3  & Poncin et al.\ (2011)~\cite{poncin2011}          & PM + MSR intersection       & IC1, IC4 & IEEE \\
V4  & Caldeira \& Cardoso (2019)                        & PM + SW process assessment  & IC1      & IEEE \\
V5  & Lemos et al.\ (2011)                              & PM for SW management        & IC1      & IEEE \\
V6  & Whittaker \& Thomason (1994)~\cite{whittaker1994} & Markov + SW testing (seminal) & IC2    & Scopus \\
V7  & Tax et al.\ (2017)~\cite{tax2017}                & PPM seminal (BPM baseline)  & IC3      & Control \\
V8  & Rogge-Solti \& Weske (2015)                       & Stochastic + PM forecasting & IC2, IC3 & Scopus \\
V9  & Rozinat et al.\ (2009)                            & PM $\rightarrow$ Simulation  & IC1, IC2 & Control \\
V10 & Gupta (2017)~\cite{gupta2017}                    & PM multi-repository SE      & IC1, IC3 & ACM \\
V11 & Jalote \& Kamma (2021)~\cite{jalote2021}         & IDE logs + Markov chains    & IC1, IC2 & Scopus \\
V12 & Ortu et al.\ (2023)~\cite{ortu2023}              & GitHub Markov + SNA         & IC2, IC4 & Scopus \\
\midrule
\multicolumn{4}{l}{Capture rate from primary queries alone} & 10/12 (83.3\%) \\
\multicolumn{4}{l}{Capture rate including snowballing}      & 12/12 (100\%) \\
\bottomrule
\end{tabular}
\end{table}

Papers V4, V5, V8, and V9 were not in the initial working set but were
recovered through snowballing during the FT phase; they are included in
the final corpus.

\section{SPMF Integration-Level Taxonomy}
\label{app:slr-spmf}

The \textit{Software Process Mining and Forecasting} (SPMF) taxonomy, used
throughout Chapter~\ref{ch:slr}, classifies studies along three orthogonal
dimensions (Table~\ref{tab:app-spmf}).

\begin{table}[htbp]
\centering
\caption{SPMF taxonomy: three dimensions with level definitions.}
\label{tab:app-spmf}
\small
\setlength{\tabcolsep}{4pt}
\begin{tabular}{l l p{7.8cm}}
\toprule
\textbf{Dimension} & \textbf{Level} & \textbf{Definition} \\
\midrule
\multirow{4}{*}{\textbf{Analytical Depth}}
  & AD1 -- Descriptive  & Discovery of as-is process models; conformance checking \\
  & AD2 -- Diagnostic   & Bottleneck detection; rework quantification; root-cause analysis \\
  & AD3 -- Predictive   & Outcome prediction; remaining-time estimation \\
  & AD4 -- Forecasting  & Distributional forecasts; probabilistic completion estimates \\
\midrule
\multirow{5}{*}{\textbf{SDLC Coverage}}
  & SC-R & Requirements and design \\
  & SC-D & Development (coding, commits, pull requests) \\
  & SC-Q & Quality (testing, reviews, defects) \\
  & SC-C & CI/CD and delivery \\
  & SC-O & Operations and maintenance \\
\midrule
\multirow{4}{*}{\textbf{Technique Integration}}
  & TI-L0 & Single technique (PM only, ML only, or stochastic only) \\
  & TI-L1 & Two techniques present in same study without explicit pipeline contract \\
  & TI-L2 & Partial bridge: two techniques linked, but incomplete or missing ML closure \\
  & TI-L3 & Unified pipeline: PM $\rightarrow$ Markov $\rightarrow$ MC $\rightarrow$ ML
             with formal inter-stage contracts \\
\bottomrule
\end{tabular}
\end{table}

The L0--L3 scale in dimension TI maps to the integration-level coding used
in Sections~\ref{sec:slr-rq3} and~\ref{app:slr-top30}. PATHCAST is positioned
at TI-L3, AD4, with SC coverage spanning SC-D, SC-Q, and SC-C.


% =============================================================================
% APPENDIX E — Reading List and Pipeline Artefacts
% =============================================================================
\chapter{Reading List and Pipeline Artefacts}
\label{app:E}

This appendix presents the ranked reading list of 32 primary studies
selected for in-depth synthesis and documents the scripts and persistent
artefacts of the review pipeline.

\section{Multi-Criteria Scoring Rubric}
\label{app:slr-ranking-criteria}

Table~\ref{tab:app-ranking-criteria} presents the 15-point rubric used to
rank studies. Three seminal papers (Cook 1995, Cook 1998, Rubin 2007) are
included regardless of score.

\begin{table}[htbp]
\centering
\caption{Multi-criteria scoring rubric for the top-32 reading list.}
\label{tab:app-ranking-criteria}
\small
\setlength{\tabcolsep}{4pt}
\begin{tabular}{p{7.5cm} c}
\toprule
\textbf{Criterion} & \textbf{Points} \\
\midrule
IC2 present (stochastic modelling in SE)          & +3 \\
IC3 present (forecasting of process metrics)      & +3 \\
IC1 present (process mining in SE)                & +2 \\
IC4 present (repository mining for process)       & +1 \\
PDF locally available                             & +2 \\
Published in 2020 or later                        & +1 \\
High-impact venue (IEEE/ACM/IST/JSS/IS/TSE)      & +1 \\
Empirical study type (case\_study or experiment)  & +1 \\
Contribution type = prediction or simulation      & +1 \\
\midrule
\textbf{Maximum score}                            & \textbf{15} \\
\bottomrule
\end{tabular}
\end{table}

\section{Top-32 Ranked Reading List}
\label{app:slr-top30}

\begin{table}[p]
\centering
\caption{Top-32 reading list with score, IC cluster, and main finding
  summary. Papers marked $\dagger$ are forced-inclusion seminal works.}
\label{tab:app-top30}
\small
\setlength{\tabcolsep}{3pt}
\begin{tabular}{c c p{3.8cm} c c p{5.7cm}}
\toprule
\textbf{Rank} & \textbf{Pts} & \textbf{Authors (year)} &
  \textbf{ICs} & \textbf{PDF} & \textbf{Main finding (summary)} \\
\midrule
 1 & 14 & Joshi \& Kahani (2024)~\cite{joshi2024}
         & IC2,IC3,IC4 & \checkmark
         & RL/MDP for GitHub PR outcome prediction \\
 2 & 14 & Li et al.\ (2020)~\cite{li2020}
         & IC2,IC3 & \checkmark
         & Hybrid SD+DES simulation for industrial project forecasting \\
 3 & 13 & Guinea-Cabrera et al.\ (2025)~\cite{guinea2025}
         & IC1,IC2,IC3 & \checkmark
         & PRIMAD Digital Twin for SDLC (pm4py + Akka + LightGBM), L2 \\
 4 & 13 & Ortu et al.\ (2023)~\cite{ortu2023}
         & IC2,IC4 & \checkmark
         & Fault-insertion/fixing Markov + SNA in Apache projects \\
 5 & 13 & Ben Mesmia et al.\ (2021)~\cite{benmesmia2021}
         & IC2,IC3 & \checkmark
         & Non-Markovian SPN for DevOps workflow duration prediction \\
 6 & 12 & Buliga et al.\ (2025)~\cite{buliga2025}
         & IC1,IC3 & \checkmark
         & Generative AI + business process simulation for what-if \\
 7 & 12 & Bhadra et al.\ (2023)~\cite{bhadra2023}
         & IC2,IC3 & \checkmark
         & SPN model of CI/CD with reliability sensitivity analysis \\
 8 & 12 & Bhadra (2022)~\cite{bhadra2022}
         & IC2,IC3 & \checkmark
         & SPN model of CI/CD pipeline delivery probability \\
 9 & 12 & Massitela et al.\ (2018)~\cite{massitela2018}
         & IC2,IC3 & \checkmark
         & Stochastic Automata Network for Waterfall team estimation \\
10 & 11 & Incerto et al.\ (2025)~\cite{incerto2025}
         & IC1,IC2 & \checkmark
         & VLMC stochastic conformance; outperforms 18 state-of-art techniques \\
11 & 11 & Jo \& Kwon (2024)~\cite{jo2024}
         & IC2,IC4 & \checkmark
         & Probabilistic model checking (CTL) on GitHub repos \\
12 & 11 & Jo et al.\ (2023)~\cite{jo2023}
         & IC2,IC4 & \checkmark
         & DTMC collaboration patterns in GitHub (5 repos) \\
13 & 11 & L\'opez-Pintado \& Dumas (2023)~\cite{lopezpintado2023}
         & IC1,IC2 & ---
         & Probabilistic resource calendars discovered from event logs \\
14 & 11 & Caldeira et al.\ (2022)~\cite{caldeira2022}
         & IC1,IC3 & \checkmark
         & Process complexity $\to$ cyclomatic complexity ($r \approx 0.43$) \\
15 & 11 & Jalote \& Kamma (2021)~\cite{jalote2021}
         & IC1,IC2 & \checkmark
         & Markov chains on IDE logs for programmer productivity \\
16 & 11 & Lunesu et al.\ (2021)~\cite{lunesu2021}
         & IC2,IC3 & \checkmark
         & Monte Carlo risk assessment for agile software development \\
17 & 11 & Nafreen et al.\ (2020)~\cite{nafreen2020}
         & IC2,IC3 & ---
         & SRGM + defect tracking for hybrid reliability forecasting \\
18 & 11 & Gupta (2017)~\cite{gupta2017}
         & IC1,IC3,IC4 & ---
         & Inductive miner + predictive analytics for SW maintenance \\
19 & 10 & Salmani et al.\ (2022)
         & IC1,IC3 & \checkmark
         & Process mining for requirements engineering in MDPs \\
20 & 10 & Jeon et al.\ (2015)~\cite{jeon2015}
         & IC2,IC3,IC4 & ---
         & Markov-based probabilistic risk prediction from repo data \\
21 & 10 & Tian \& Noore (2005)~\cite{tian2005}
         & IC2,IC3 & ---
         & Markov defect-removal modelling under code churn \\
22 & 10 & Cook \& Wolf (1995)$^\dagger$~\cite{cook1995}
         & IC1,IC2 & \checkmark
         & First automated process discovery from event traces \\
23 &  9 & Pourbafrani et al.\ (2021)~\cite{pourbafrani2021}
         & IC1,IC3 & \checkmark
         & SIMPT: process tree + interactive what-if simulation tool \\
24 &  9 & Tyagi et al.\ (2021)~\cite{tyagi2021}
         & IC2,IC3 & ---
         & State-space Markov framework for agile planning and tracking \\
25 &  9 & Razi et al.\ (2019)
         & IC1,IC2 & \checkmark
         & Conformance checking for SDLC process compliance \\
26 &  9 & Krismayer et al.\ (2019)
         & IC1,IC4 & \checkmark
         & Declarative PM for repository-based process monitoring \\
27 &  9 & Baskoro et al.\ (2019)
         & IC3,IC4 & ---
         & ML-based effort prediction from repository data \\
28 &  9 & (Anon.\ 2017)
         & IC2,IC3 & ---
         & Probabilistic defect prediction for software product lines \\
29 &  9 & Magennis (2015)~\cite{magennis2015}
         & IC2,IC3 & ---
         & Monte Carlo cycle-time economic analysis \\
30 &  9 & Washizaki et al.\ (2015)~\cite{washizaki2015}
         & IC2,IC3 & ---
         & GSRM release-time prediction for open-source projects \\
31 &  9 & Cook \& Wolf (1998)$^\dagger$~\cite{cook1998}
         & IC1,IC2 & ---
         & Markov/FSA/NN process discovery from event-based data \\
32 &  2 & Rubin et al.\ (2007)$^\dagger$~\cite{rubin2007}
         & IC1 & \checkmark
         & First named PM framework for SDLC (ProM + SCM logs) \\
\bottomrule
\end{tabular}
\end{table}

\section{Review Pipeline Scripts}
\label{app:slr-pipeline}

The review was conducted using a custom Python pipeline (\texttt{pipeline/})
operating on a structured CSV data store. Table~\ref{tab:app-pipeline}
documents each script with its function and main invocation command.

\begin{table}[htbp]
\centering
\caption{Review pipeline scripts and their roles.}
\label{tab:app-pipeline}
\small
\setlength{\tabcolsep}{3pt}
\begin{tabular}{p{4.2cm} p{4.5cm} p{4.0cm}}
\toprule
\textbf{Script} & \textbf{Function} & \textbf{Main command} \\
\midrule
\texttt{pipeline/dedup.py} &
  Cross-source deduplication (DOI + fuzzy title) &
  \texttt{python main.py dedup} \\[2pt]
\texttt{pipeline/enrich.py} &
  Abstract enrichment cascade (4 sources) &
  \texttt{python main.py enrich} \\[2pt]
\texttt{pipeline/screening.py} &
  T/A LLM screening via Batches API &
  \texttt{python main.py screen -{}-run} \\[2pt]
\texttt{pipeline/fulltext.py} &
  FT LLM screening + manual review &
  \texttt{python main.py fulltext -{}-run} \\[2pt]
\texttt{pipeline/pdf\_downloader.py} &
  PDF retrieval (Unpaywall, S2, OA) &
  \texttt{python main.py pdf} \\[2pt]
\texttt{pipeline/extract\_prep.py} &
  S2 metadata enrichment + PDF copy &
  \texttt{python pipeline/} \texttt{extract\_prep.py} \\[2pt]
\texttt{pipeline/extract\_llm.py} &
  Structured data extraction (LLM, 169 papers) &
  \texttt{python pipeline/} \texttt{extract\_llm.py -{}-run} \\[2pt]
\texttt{pipeline/synth\_llm.py} &
  Synthesis paragraph generation per IC cluster &
  \texttt{python pipeline/} \texttt{synth\_llm.py} \\[2pt]
\texttt{pipeline/finalization.py} &
  PRISMA snapshot, EC5 closure, artefact export &
  \texttt{python main.py finalize} \\
\bottomrule
\end{tabular}
\end{table}

\section{Persistent Artefact Registry}
\label{app:slr-artefacts}

Table~\ref{tab:app-artefacts} lists the persistent artefacts produced by the
review, with their role, location, and current state.

\begin{table}[htbp]
\centering
\caption{Persistent artefacts of the review.}
\label{tab:app-artefacts}
\small
\setlength{\tabcolsep}{3pt}
\begin{tabular}{p{6.5cm} p{3.0cm} p{2.8cm}}
\toprule
\textbf{File} & \textbf{Role} & \textbf{State} \\
\midrule
\texttt{results/screening/} \texttt{ft\_screening\_results.csv}
  & 886-paper decision log & Final \\[3pt]
\texttt{results/extraction/} \texttt{extraction\_template.csv}
  & 169-paper extraction data & 169/169 \\[3pt]
\texttt{results/final\_review/} \texttt{included\_studies\_current.csv}
  & Included study set & 169 rows \\[3pt]
\texttt{results/final\_review/} \texttt{top30\_reading\_list.csv}
  & Ranked reading list & 32 rows \\[3pt]
\texttt{results/final\_review/} \texttt{cap3\_synthesis\_v2.tex}
  & LLM synthesis paragraphs & For review \\[3pt]
\texttt{results/final\_review/} \texttt{missing\_references.bib}
  & BibTeX for 32 missing keys & Ready to merge \\[3pt]
\texttt{results/final\_review/} \texttt{pending\_inaccessible\_closed.csv}
  & EC5 closure record & 148 rows \\[3pt]
\texttt{results/final\_review/} \texttt{prisma\_summary\_current.txt}
  & PRISMA snapshot & Final \\[3pt]
\texttt{results/} \texttt{SLR\_EXECUCAO\_ATE\_AGORA.md}
  & Execution log (all sessions) & Up to date \\
\bottomrule
\end{tabular}
\end{table}

% The \appendix command causes subsequent \chapter{} calls to be labelled
% as APPENDIX A, B, C, etc. per ABNT NBR 14724.
% =============================================================================


% =============================================================================
% APPENDIX A — Search Protocol: Search Strings
% =============================================================================
\chapter{Search Protocol: Search Strings}
\label{app:A}

This appendix presents the four search strings executed across the five
digital libraries during the search phase of the PATHCAST SLR. Each string
was adapted to the syntax and field conventions of the target database; the
block structure (P and I) is preserved across all adaptations. Strings S1--S3
correspond to the principal, complementary (MSR), and frontier (stochastic
$\wedge$ PM) blocks, respectively; S4 is a supplementary string executed in
Scopus only, targeting Markov chain models in software testing.

\section{Search String Summary by Database}
\label{app:slr-strings}

\begin{table}[htbp]
\centering
\caption{Search strings by database and field scope.}
\label{tab:app-strings}
\small
\setlength{\tabcolsep}{4pt}
\begin{tabular}{p{2.4cm} c p{9.0cm}}
\toprule
\textbf{Database} & \textbf{Strings} & \textbf{Field / syntax notes} \\
\midrule
Scopus       & S1, S2, S3, S4 & \texttt{TITLE-ABS-KEY(...)};
                                  \texttt{DOCTYPE(ar OR cp)};
                                  \texttt{LANGUAGE(english)};
                                  1994--2026 \\
IEEE Xplore  & S1 (split)     & All Metadata; Boolean AND of sub-blocks;
                                  Conference + Journals; 1994--2026 \\
Web of Science & S1, S2, S3   & Topic \texttt{TS=(...)}; Article, Proceedings;
                                  English; 1994--2026 \\
SpringerLink & S1 (simplified)& Title + Abstract; 8 separate exports;
                                  Chapter + Article; English \\
ACM DL       & S1, S2         & Abstract; 2 queries; date 1994--2026 \\
\bottomrule
\end{tabular}
\end{table}

\section{String S1 — Principal ($\text{P} \wedge \text{I}$)}
\label{app:slr-s1}

Block P (software engineering context):

\begin{quote}
\small\ttfamily
("software engineering" OR "software development" OR "software process"
OR "software lifecycle" OR "software life cycle" OR "SDLC" OR "development
lifecycle" OR "development life cycle" OR "DevOps" OR "CI/CD" OR
"continuous integration" OR "continuous delivery" OR "agile development"
OR "agile process" OR "issue tracking" OR "bug tracking" OR "version
control" OR "code review" OR "pull request" OR "software project"
OR "software repository" OR "software maintenance" OR "software
evolution")
\end{quote}

Block I (techniques):

\begin{quote}
\small\ttfamily
("process mining" OR "process discovery" OR "conformance checking" OR
"workflow mining" OR "predictive process monitoring" OR "process prediction"
OR "remaining time prediction" OR "process forecasting" OR "event log" OR
"process simulation" OR "Monte Carlo simulation" OR "Markov chain" OR
"Markov model" OR "stochastic process model" OR "transition probability"
OR "absorbing chain" OR "lead time prediction" OR "cycle time prediction")
\end{quote}

\section{String S2 — Complementary (MSR $\rightarrow$ Process Modeling)}
\label{app:slr-s2}

\begin{quote}
\small\ttfamily
("mining software repositories" OR "software repository mining" OR
"GitHub mining" OR "Jira mining") AND ("process model" OR "Markov" OR
"Monte Carlo" OR "stochastic" OR "lead time" OR "cycle time" OR "throughput"
OR "transition matrix")
\end{quote}

\section{String S3 — Frontier (Stochastic $\wedge$ PM $\wedge$ SE)}
\label{app:slr-s3}

\begin{quote}
\small\ttfamily
("Monte Carlo" OR "Markov chain" OR "stochastic simulation" OR
"probabilistic forecast" OR "transition matrix") AND ("process mining" OR
"event log" OR "workflow mining" OR "process model" OR "business process")
AND ("software development" OR "software engineering" OR "software process"
OR "DevOps" OR "SDLC" OR "issue tracker" OR "commit log" OR "build pipeline")
\end{quote}

\section{String S4 — Markov Chain Software Testing (Scopus supplement)}
\label{app:slr-s4}

\begin{quote}
\small\ttfamily
("Markov chain" OR "Markov model" OR "hidden Markov" OR "stochastic automata")
AND ("software testing" OR "test coverage" OR "software reliability" OR
"defect prediction" OR "fault prediction")
\end{quote}


% =============================================================================
% APPENDIX B — LLM Screening Prompts
% =============================================================================
\chapter{LLM Screening Prompts}
\label{app:B}

The review employed \textit{claude-haiku-4-5-20251001} as primary screener in
two phases: title/abstract (T/A) and full-text (FT). Each phase uses a
\emph{system prompt} that encodes the scope, criteria, and output schema,
paired with a per-paper \emph{user prompt} that supplies the bibliographic
data. All prompts are version-controlled in
\texttt{config/screening\_criteria.py} in the review repository.

\section{Phase 1: Title/Abstract Screening}
\label{app:slr-ta-prompt}

\subsection*{System Prompt (T/A Phase)}

\begin{quote}
\small\ttfamily
You are an expert in systematic literature review (SLR) in the area of
process mining and software engineering. Your task is to perform the
title/abstract (T/A) screening for SLR PATHCAST.

\textbf{SLR PATHCAST scope:}
Intersection of (1) process mining, (2) stochastic modelling,
(3) forecasting/prediction, and (4) software development processes.
Focus: studies that analyse, discover, or predict software process
characteristics using event logs, repositories, or stochastic models.

\textbf{INCLUSION CRITERIA} (at least ONE must be satisfied to include):
IC1 -- Process Mining in Software: applies PM (discovery, conformance,
workflow mining, event log analysis, PPM) to SW artefacts (commits,
issues, PRs, CI/CD).
IC2 -- Stochastic Modelling of SW Processes: uses Markov chains, Monte
Carlo, stochastic Petri nets, or transition matrices to model SW dev
processes.
IC3 -- Process Metric Forecasting: predicts lead time, cycle time,
remaining time, throughput, or defect rates in SW processes using
data or event logs.
IC4 -- Repository Mining for Process: mines repositories (GitHub, Jira,
VCS) to discover or improve SW development process models.

\textbf{EXCLUSION CRITERIA} (any ONE is sufficient to exclude):
EC1 -- Non-SW domain: applied to healthcare, manufacturing, finance,
etc., with no link to SW processes.
EC2 -- SW as tool: "software" is the implementation platform, not the
process under study.
EC3 -- Algorithm without SW application: purely theoretical method,
no evaluation in SW processes.
EC4 -- Out-of-scope secondary study: survey/SLR not specifically
addressing PM or stochastic methods in SE.

\textbf{UNCERTAIN CASES (maybe):}
Relevant title without abstract; generic BPM article with possible SW
application; ambiguous terminology.

\textbf{STRUCTURED FIELDS} (use only the allowed enums):
evidence\_tags, software\_context, stochastic\_method,
forecast\_target, process\_data\_source, confidence.

Respond ONLY with valid JSON, no additional text.
\end{quote}

\subsection*{User Prompt Template (T/A Phase)}

\begin{quote}
\small\ttfamily
Evaluate the paper below for SLR PATHCAST.

\textbf{Title:} \{title\}
\textbf{Abstract:} \{abstract\}
\textbf{Venue/Type:} \{venue\} | \{doc\_type\} | \{year\}
\textbf{Source database:} \{source\_db\}
\textbf{Abstract origin:} \{abstract\_source\}

Respond with JSON in the exact format:
\{
  "decision": "include" | "exclude" | "maybe",
  "rationale": "<1-2 sentences justifying the decision>",
  "matched\_ic": ["IC1", "IC2"],
  "matched\_ec": ["EC1"],
  "evidence\_tags": ["process\_mining", "forecasting"],
  "software\_context": "software\_development\_process",
  "stochastic\_method": "markov\_chain",
  "forecast\_target": "remaining\_time",
  "process\_data\_source": "event\_logs",
  "confidence": "medium"
\}

Rules:
- "include" if $\geq$1 IC met and no decisive EC
- "exclude" if an EC clearly applies, or paper clearly meets no IC
- "maybe" if genuine doubt (title without abstract, domain ambiguity)
- matched\_ic and matched\_ec must be [] when not applicable
\end{quote}

\section{Phase 2: Full-Text Screening}
\label{app:slr-ft-prompt}

\subsection*{System Prompt (FT Phase)}

The FT system prompt is identical to the T/A system prompt above, with
three additional instructions that encode the stricter standard of the
full-text phase:

\begin{quote}
\small\ttfamily
THIS IS THE FULL-TEXT PHASE -- apply criteria rigorously:
- If the abstract clearly confirms $\geq$1 IC with no decisive EC
  $\rightarrow$ "include"
- If the abstract clearly confirms no IC applies, or an EC applies
  $\rightarrow$ "exclude"
- Use "pending" ONLY if the abstract is genuinely insufficient for
  a decision (absent or extremely short) AND the title suggests
  marginal relevance

When in doubt between include and exclude, prefer exclude -- this
phase requires confidence.

YOU MUST NEVER invent, complete, or infer a missing abstract.
\end{quote}

\subsection*{User Prompt Template (FT Phase)}

The FT user prompt extends the T/A template with a \textbf{T/A context block}
that supplies the previous-phase decision, rationale, and matched ICs:

\begin{quote}
\small\ttfamily
Evaluate the paper below for SLR PATHCAST -- FULL-TEXT screening (phase 2).

\textbf{Title:} \{title\}
\textbf{Abstract:} \{abstract\}
\textbf{Venue/Type:} \{venue\} | \{doc\_type\} | \{year\}
\textbf{Source database:} \{source\_db\}
\textbf{Abstract origin:} \{abstract\_source\}

\textbf{T/A phase context:}
  - Previous decision: \{ta\_decision\}
  - Previous rationale: \{ta\_rationale\}
  - ICs identified in phase 1: \{ta\_matched\_ic\}

Re-evaluate rigorously for the full-text phase.
Respond with JSON in the exact format:
\{
  "decision": "include" | "exclude" | "pending",
  [... same fields as T/A phase ...]
\}
\end{quote}

\section{Phase 2b: Screen-Blanks Pass}
\label{app:slr-blanks-prompt}

A third pass targeted the 186 papers without a decision after both prior
phases (primarily Band~E papers with no available abstract). The screen-blanks
pass reuses the same FT system and user prompts, with the abstract field
populated by the placeholder ``Abstract not available'' and the T/A context
block showing \texttt{ta\_decision~=~maybe} and
\texttt{ta\_matched\_ic~=~none}. Papers for which the LLM returned
\texttt{pending} and whose title analysis suggested low relevance were
subsequently closed under EC5 (full text inaccessible; $n = 148$).

\section{Model Parameters per Screening Phase}
\label{app:slr-prompt-params}

Table~\ref{tab:app-prompt-params} summarises the model parameters used
across all three screening phases and the data extraction step.

\begin{table}[htbp]
\centering
\caption{LLM parameters used at each screening phase.}
\label{tab:app-prompt-params}
\small
\setlength{\tabcolsep}{4pt}
\begin{tabular}{l l r r l}
\toprule
\textbf{Phase} & \textbf{Model} & \textbf{Max tok.} &
  \textbf{Batch} & \textbf{API} \\
\midrule
T/A screening      & claude-haiku-4-5-20251001 & 512  & 500 & Anthropic Batches \\
FT screening       & claude-haiku-4-5-20251001 & 512  & 100 & Anthropic Batches \\
Screen-blanks pass & claude-haiku-4-5-20251001 & 512  & 186 & Anthropic Batches \\
Data extraction    & claude-haiku-4-5-20251001 & 1024 & 100 & Anthropic Batches \\
Synthesis (ch.~3)  & claude-sonnet-4-6         & 2048 & --- & Anthropic Messages \\
\bottomrule
\end{tabular}
\end{table}


% =============================================================================
% APPENDIX C — Data Extraction Form
% =============================================================================
\chapter{Data Extraction Form}
\label{app:C}

Data extraction for all 169 included studies was performed with
\textit{claude-haiku-4-5-20251001} via the Anthropic Batches API. Papers
with a locally available PDF had their text pre-extracted with
\texttt{pdfplumber} (first 8 pages, maximum 6,000 characters). Papers
without a PDF used the enriched abstract from Semantic Scholar.
The prompt below was submitted once per paper.

\section{Extraction Prompt}
\label{app:slr-extraction-prompt}

\begin{quote}
\small\ttfamily
You are a researcher performing data extraction for SLR PATHCAST on
Process Mining and Stochastic Modelling in Software Processes.

\textbf{Title:} \{title\}
\textbf{Authors:} \{authors\}
\textbf{Year:} \{year\} | \textbf{Venue:} \{venue\}
\textbf{Inclusion criteria met:} \{ics\}
\{content\_block\}

Return ONLY a valid JSON object with exactly these fields:
\{
  "research\_question": "main research question (1-2 sentences)",
  "study\_type": "case\_study | experiment | survey | tool |
                 theoretical | simulation | mixed",
  "research\_contribution": "discovery | conformance | enhancement |
       prediction | simulation | framework | hybrid",
  "pm\_technique": "alpha | inductive\_miner | heuristic |
       conformance\_checking | social\_network | declarative | other | none",
  "stochastic\_technique": "markov\_chain | stochastic\_petri\_net |
       monte\_carlo | system\_dynamics | other | none",
  "software\_artifact": "commits | issues | ci\_cd | ide\_logs |
       vcs | jira | github | other | none",
  "software\_process": "development | testing | maintenance |
       code\_review | bug\_fixing | deployment | requirements | other",
  "dataset\_source": "open\_source | industrial | academic |
                     synthetic | not\_specified",
  "dataset\_public": "yes | no | partial | not\_mentioned",
  "tool\_used": "tool names (e.g.\ ProM, pm4py, Disco) or not\_mentioned",
  "main\_finding": "main result/contribution (2-3 sentences)",
  "limitations": "reported limitations (1-2 sentences) or not\_mentioned",
  "replication\_package": "yes | no | partial | not\_mentioned"
\}

Respond ONLY with JSON, no additional text.
\end{quote}

\section{Extraction Form Fields}
\label{app:slr-extraction-fields}

Table~\ref{tab:app-extraction-fields} describes all fields of the structured
extraction form stored in \texttt{extraction\_template.csv}.

\begin{table}[htbp]
\centering
\caption{Data extraction form: all fields in \texttt{extraction\_template.csv}.}
\label{tab:app-extraction-fields}
\footnotesize
\setlength{\tabcolsep}{3pt}
\begin{tabular}{p{3.2cm} p{4.4cm} p{4.6cm}}
\toprule
\textbf{Field} & \textbf{Description} & \textbf{Allowed values / format} \\
\midrule
\multicolumn{3}{l}{\textit{Identification}} \\
\texttt{internal\_id}   & 8-char hash (dedup key)      & --- \\
\texttt{title}          & Paper title                  & free text \\
\texttt{doi}            & Digital Object Identifier    & URI \\
\texttt{year}           & Publication year             & YYYY \\
\texttt{source\_db}     & Source database              & scopus, ieee, acm, springer, wos, snowball \\
\texttt{ft\_matched\_ic}& Inclusion criteria met       & IC1$|$IC2$|$IC3$|$IC4 \\
\texttt{ft\_screened\_by}& Who made FT decision        & llm, manual, llm\_pdf \\
\midrule
\multicolumn{3}{l}{\textit{Enriched metadata (Semantic Scholar)}} \\
\texttt{authors}        & Author list                  & semicolon-separated \\
\texttt{venue}          & Publication venue            & free text \\
\texttt{journal\_name}  & Journal/conference name      & free text \\
\texttt{abstract}       & Paper abstract               & free text \\
\texttt{pdf\_available} & PDF locally present          & yes / no \\
\midrule
\multicolumn{3}{l}{\textit{LLM-extracted fields}} \\
\texttt{research\_question} & Main RQ               & 1--2 sentences \\
\texttt{study\_type}        & Study design          & case\_study $|$ experiment $|$ survey $|$ tool $|$ theoretical $|$ simulation $|$ mixed \\
\texttt{research\_contribution} & Contribution type & discovery $|$ conformance $|$ enhancement $|$ prediction $|$ simulation $|$ framework $|$ hybrid \\
\texttt{pm\_technique}      & PM algorithm used     & alpha $|$ inductive\_miner $|$ heuristic $|$ conformance\_checking $|$ social\_network $|$ declarative $|$ other $|$ none \\
\texttt{stochastic\_technique} & Stochastic method  & markov\_chain $|$ stochastic\_petri\_net $|$ monte\_carlo $|$ system\_dynamics $|$ other $|$ none \\
\texttt{software\_artifact} & Data source           & commits $|$ issues $|$ ci\_cd $|$ ide\_logs $|$ vcs $|$ jira $|$ github $|$ other $|$ none \\
\texttt{software\_process}  & SDLC phase            & development $|$ testing $|$ maintenance $|$ code\_review $|$ bug\_fixing $|$ deployment $|$ requirements $|$ other \\
\texttt{dataset\_source}    & Dataset provenance    & open\_source $|$ industrial $|$ academic $|$ synthetic $|$ not\_specified \\
\texttt{dataset\_public}    & Data availability     & yes $|$ no $|$ partial $|$ not\_mentioned \\
\texttt{tool\_used}         & Tools/platforms       & free text \\
\texttt{main\_finding}      & Key result            & 2--3 sentences \\
\texttt{limitations}        & Study limitations     & 1--2 sentences \\
\texttt{replication\_package} & Artefacts available & yes $|$ no $|$ partial $|$ not\_mentioned \\
\midrule
\multicolumn{3}{l}{\textit{Manual fields (reading phase)}} \\
\texttt{quality\_score}     & QA checklist score (0--8) & numeric \\
\texttt{extraction\_notes}  & Free-form reviewer note   & free text \\
\bottomrule
\end{tabular}
\end{table}

\section{Structured Output Controlled Vocabulary}
\label{app:slr-schema}

Table~\ref{tab:app-enum-schema} defines the complete controlled vocabulary
for the structured fields returned by the LLM at both screening phases.
These enums are version-controlled in \texttt{config/screening\_criteria.py}
and were not changed between screening rounds to ensure consistency.

\begin{table}[htbp]
\centering
\caption{Controlled vocabulary (enums) for LLM structured output.}
\label{tab:app-enum-schema}
\small
\setlength{\tabcolsep}{4pt}
\begin{tabular}{p{3.2cm} p{9.4cm}}
\toprule
\textbf{Field} & \textbf{Allowed values} \\
\midrule
\texttt{decision} (T/A) &
  \texttt{include}, \texttt{exclude}, \texttt{maybe} \\
\texttt{decision} (FT) &
  \texttt{include}, \texttt{exclude}, \texttt{pending} \\
\texttt{software\_context} &
  \texttt{software\_development\_process}, \texttt{repository\_mining},
  \texttt{ci\_cd}, \texttt{issue\_bug\_workflow}, \texttt{software\_testing},
  \texttt{requirements\_engineering}, \texttt{software\_project\_management},
  \texttt{unclear}, \texttt{not\_software\_process} \\
\texttt{stochastic\_method} &
  \texttt{none}, \texttt{markov\_chain}, \texttt{hidden\_markov\_model},
  \texttt{monte\_carlo}, \texttt{stochastic\_petri\_net},
  \texttt{bayesian\_model}, \texttt{probabilistic\_model},
  \texttt{simulation}, \texttt{queueing\_model},
  \texttt{other\_stochastic}, \texttt{unclear} \\
\texttt{forecast\_target} &
  \texttt{none}, \texttt{lead\_time}, \texttt{cycle\_time},
  \texttt{remaining\_time}, \texttt{throughput}, \texttt{defect\_rate},
  \texttt{build\_outcome}, \texttt{reliability},
  \texttt{completion\_time}, \texttt{other\_process\_metric},
  \texttt{unclear} \\
\texttt{process\_data\_source} &
  \texttt{none}, \texttt{event\_logs}, \texttt{version\_control},
  \texttt{issue\_tracker}, \texttt{pull\_requests}, \texttt{ci\_cd\_logs},
  \texttt{software\_repository\_mixed}, \texttt{synthetic\_data},
  \texttt{simulated\_process}, \texttt{survey\_or\_secondary},
  \texttt{unclear} \\
\texttt{confidence} &
  \texttt{low}, \texttt{medium}, \texttt{high} \\
\texttt{evidence\_tags} &
  \texttt{process\_mining}, \texttt{software\_process},
  \texttt{repository\_mining}, \texttt{stochastic\_modeling},
  \texttt{forecasting}, \texttt{event\_log}, \texttt{version\_control},
  \texttt{issue\_tracking}, \texttt{pull\_requests}, \texttt{ci\_cd},
  \texttt{markov}, \texttt{hidden\_markov\_model}, \texttt{monte\_carlo},
  \texttt{stochastic\_petri\_net}, \texttt{bayesian\_model},
  \texttt{simulation}, \texttt{lead\_time}, \texttt{cycle\_time},
  \texttt{remaining\_time}, \texttt{throughput}, \texttt{defect\_prediction},
  \texttt{build\_prediction}, \texttt{reliability},
  \texttt{insufficient\_abstract}, \texttt{full\_text\_required} \\
\bottomrule
\end{tabular}
\end{table}


% =============================================================================
% APPENDIX D — Assessment and Classification Instruments
% =============================================================================
\chapter{Assessment and Classification Instruments}
\label{app:D}

This appendix gathers the instruments used to classify, enrich, and evaluate
papers throughout the screening and selection process: the priority band
taxonomy, the abstract enrichment cascade, the quality assessment checklist,
the control paper set, and the SPMF integration-level taxonomy.

\section{Priority Bands for the Full-Text Review Queue}
\label{app:slr-bands}

Prior to full-text screening, the 886-paper queue was stratified into five
priority bands (A--E) to optimise resource allocation. Abstract enrichment
and LLM screening were applied in band order.
Table~\ref{tab:app-bands} defines the classification rules.

\begin{table}[htbp]
\centering
\caption{Priority band definitions for the full-text review queue.}
\label{tab:app-bands}
\small
\setlength{\tabcolsep}{4pt}
\begin{tabular}{c r p{8.5cm}}
\toprule
\textbf{Band} & \textbf{$n$} & \textbf{Classification rule} \\
\midrule
A &  95 & T/A decision = \texttt{include} and $\geq$2 distinct ICs matched \\
B &  20 & T/A decision = \texttt{include} and exactly 1 IC matched \\
C &  39 & T/A decision = \texttt{maybe} and abstract available \\
D & 152 & T/A decision = \texttt{maybe}, at least 1 IC signal in title, no abstract \\
E & 584 & T/A decision = \texttt{maybe}, no IC signal in title, no abstract \\
\midrule
\textbf{Total} & \textbf{886} & \\
\bottomrule
\end{tabular}
\end{table}

Band~D papers with a DOI were prioritised for manual PDF retrieval. Band~E
papers without a DOI and without an abstract were closed as EC5 (full text
inaccessible) at the end of the screen-blanks pass.

\section{Abstract Enrichment Cascade}
\label{app:slr-enrichment}

To increase the number of papers assessable by the LLM, a four-source
cascade was executed prior to full-text screening.
Algorithm~\ref{alg:enrichment} describes the procedure.

\begin{figure}[htbp]
\centering
\small
\begin{tabular}{l}
\hline
\textbf{Algorithm 1} \quad Abstract Enrichment Cascade \\
\hline
\textbf{Input:} Paper record $p$ with \texttt{doi} and/or \texttt{title} \\
\textbf{Output:} Updated $p$ with \texttt{abstract}, \texttt{authors}, \texttt{venue} \\
\hline
\textbf{1.} If $p$.\texttt{abstract} is non-empty $\rightarrow$ \textbf{return} $p$ \\
\textbf{2.} If $p$.\texttt{doi} is present: \\
\quad\textbf{2a.} Query \textbf{Semantic Scholar} Batch API (\texttt{DOI:\{doi\}}) \\
\quad\quad If match found: update \texttt{abstract}, \texttt{authors}, \texttt{venue}; \textbf{return} \\
\quad\textbf{2b.} Query \textbf{OpenAlex} (\texttt{filter=doi:\{doi\}}) \\
\quad\quad If match found: update \texttt{abstract}; \textbf{return} \\
\quad\textbf{2c.} Query \textbf{CORE} (\texttt{/search/works?q=doi:\{doi\}}) \\
\quad\quad If match found: update \texttt{abstract}; \textbf{return} \\
\quad\textbf{2d.} Query \textbf{Crossref} (\texttt{/works/\{doi\}}) \\
\quad\quad If abstract present: update; \textbf{return} \\
\textbf{3.} If $p$.\texttt{doi} is absent: \\
\quad\textbf{3a.} Query \textbf{Semantic Scholar} by title (fuzzy) \\
\quad\quad If similarity $\geq 0.85$: update; \textbf{return} \\
\textbf{4.} Mark $p$.\texttt{abstract\_source} = \texttt{missing\_or\_unverified} \\
\hline
\end{tabular}
\caption{Abstract enrichment cascade applied to all 886 full-text queue papers.}
\label{alg:enrichment}
\end{figure}

The cascade raised abstract coverage from 16.6\% (147/886) to 72.6\%
(643/886) before Round~2 screening. Of the 643 enriched abstracts,
590 were sourced from Semantic Scholar, 31 from OpenAlex, 13 from CORE,
and 9 from Crossref or the title-based fallback.
Table~\ref{tab:app-enrichment-stats} summarises the enrichment outcome
by priority band.

\begin{table}[htbp]
\centering
\caption{Abstract enrichment results by priority band.}
\label{tab:app-enrichment-stats}
\small
\setlength{\tabcolsep}{4pt}
\begin{tabular}{c r r r r}
\toprule
\textbf{Band} & \textbf{Papers} & \textbf{With abstract (pre)} &
  \textbf{Gained} & \textbf{Coverage (post)} \\
\midrule
A &  95 & 91 (95.8\%) &   4 &  95 (100.0\%) \\
B &  20 & 18 (90.0\%) &   2 &  20 (100.0\%) \\
C &  39 & 39 (100\%)  &   0 &  39 (100.0\%) \\
D & 152 &  0  (0.0\%) & 118 & 118  (77.6\%) \\
E & 584 &  0  (0.0\%) & 371 & 371  (63.5\%) \\
\midrule
\textbf{Total} & \textbf{886} & \textbf{148 (16.7\%)} &
  \textbf{495} & \textbf{643 (72.6\%)} \\
\bottomrule
\end{tabular}
\end{table}

\section{Quality Assessment Checklist}
\label{app:slr-qa}

The eight-item quality checklist (Table~\ref{tab:app-qa-full}) is applied
to each included primary study during the manual reading phase. Items
QA1--QA4 assess methodological rigour; QA5--QA7 assess empirical soundness
and reproducibility; QA8 specifically targets process model quality reporting,
which is central to PATHCAST's evaluation strategy. Studies scoring below
4/8 are excluded from the synthesis.

\begin{table}[htbp]
\centering
\caption{Quality assessment checklist with scoring guidance.}
\label{tab:app-qa-full}
\small
\setlength{\tabcolsep}{4pt}
\begin{tabular}{c p{5.5cm} p{5.5cm}}
\toprule
\textbf{ID} & \textbf{Criterion} & \textbf{Scoring guidance} \\
\midrule
QA1 & Are the research objectives clearly stated? &
  Score 1 if RQs or hypotheses are explicitly enumerated; 0 otherwise. \\
QA2 & Is the software engineering context described in sufficient detail? &
  Score 1 if the SDLC phase, artefact type, and process scope are defined. \\
QA3 & Is the data source (event log or repository) described reproducibly? &
  Score 1 if dataset origin, collection method, and key statistics are given. \\
QA4 & Is the PM or stochastic technique formally defined? &
  Score 1 if algorithm, parameters, and configuration are specified. \\
QA5 & Are results validated empirically? &
  Score 1 if validation uses real or semi-real data; 0 for simulation-only. \\
QA6 & Are threats to validity discussed? &
  Score 1 if at least internal and external validity are addressed. \\
QA7 & Is the study reproducible (data and/or code available)? &
  Score 1 if a replication package or open dataset is linked or described. \\
QA8 & Are process model quality metrics reported (fitness, precision)? &
  Score 1 if at least one formal quality metric is quantified. \\
\midrule
\multicolumn{2}{l}{\textbf{Total score range}} & 0--8 (threshold: $\geq$4) \\
\bottomrule
\end{tabular}
\end{table}

Current QA status: 76 of 169 studies assessed (full Round~1 set); mean
score = 5.18 (SD = 1.13), median = 5.0 (IQR: 4--6); 5 studies scored
below the threshold and were excluded from the synthesis. The remaining
93 studies (Rounds~2+) are pending manual QA assessment.

\section{Control Paper Set}
\label{app:slr-control}

The 12-paper control set was used to validate the search strings before
execution. Table~\ref{tab:app-control} lists all control papers with their
IC profile and the source through which each was captured.

\begin{table}[htbp]
\centering
\caption{Complete control paper set for search string validation.}
\label{tab:app-control}
\small
\setlength{\tabcolsep}{3pt}
\begin{tabular}{c p{4.0cm} p{3.8cm} l l}
\toprule
\textbf{\#} & \textbf{Paper} & \textbf{Justification} & \textbf{ICs} &
  \textbf{Captured via} \\
\midrule
V1  & Cook \& Wolf (1998)~\cite{cook1998}              & Seminal PM + SE             & IC1, IC2 & ACM \\
V2  & Rubin et al.\ (2007)~\cite{rubin2007}            & First named PM framework for SW & IC1 & Scopus \\
V3  & Poncin et al.\ (2011)~\cite{poncin2011}          & PM + MSR intersection       & IC1, IC4 & IEEE \\
V4  & Caldeira \& Cardoso (2019)                        & PM + SW process assessment  & IC1      & IEEE \\
V5  & Lemos et al.\ (2011)                              & PM for SW management        & IC1      & IEEE \\
V6  & Whittaker \& Thomason (1994)~\cite{whittaker1994} & Markov + SW testing (seminal) & IC2    & Scopus \\
V7  & Tax et al.\ (2017)~\cite{tax2017}                & PPM seminal (BPM baseline)  & IC3      & Control \\
V8  & Rogge-Solti \& Weske (2015)                       & Stochastic + PM forecasting & IC2, IC3 & Scopus \\
V9  & Rozinat et al.\ (2009)                            & PM $\rightarrow$ Simulation  & IC1, IC2 & Control \\
V10 & Gupta (2017)~\cite{gupta2017}                    & PM multi-repository SE      & IC1, IC3 & ACM \\
V11 & Jalote \& Kamma (2021)~\cite{jalote2021}         & IDE logs + Markov chains    & IC1, IC2 & Scopus \\
V12 & Ortu et al.\ (2023)~\cite{ortu2023}              & GitHub Markov + SNA         & IC2, IC4 & Scopus \\
\midrule
\multicolumn{4}{l}{Capture rate from primary queries alone} & 10/12 (83.3\%) \\
\multicolumn{4}{l}{Capture rate including snowballing}      & 12/12 (100\%) \\
\bottomrule
\end{tabular}
\end{table}

Papers V4, V5, V8, and V9 were not in the initial working set but were
recovered through snowballing during the FT phase; they are included in
the final corpus.

\section{SPMF Integration-Level Taxonomy}
\label{app:slr-spmf}

The \textit{Software Process Mining and Forecasting} (SPMF) taxonomy, used
throughout Chapter~\ref{ch:slr}, classifies studies along three orthogonal
dimensions (Table~\ref{tab:app-spmf}).

\begin{table}[htbp]
\centering
\caption{SPMF taxonomy: three dimensions with level definitions.}
\label{tab:app-spmf}
\small
\setlength{\tabcolsep}{4pt}
\begin{tabular}{l l p{7.8cm}}
\toprule
\textbf{Dimension} & \textbf{Level} & \textbf{Definition} \\
\midrule
\multirow{4}{*}{\textbf{Analytical Depth}}
  & AD1 -- Descriptive  & Discovery of as-is process models; conformance checking \\
  & AD2 -- Diagnostic   & Bottleneck detection; rework quantification; root-cause analysis \\
  & AD3 -- Predictive   & Outcome prediction; remaining-time estimation \\
  & AD4 -- Forecasting  & Distributional forecasts; probabilistic completion estimates \\
\midrule
\multirow{5}{*}{\textbf{SDLC Coverage}}
  & SC-R & Requirements and design \\
  & SC-D & Development (coding, commits, pull requests) \\
  & SC-Q & Quality (testing, reviews, defects) \\
  & SC-C & CI/CD and delivery \\
  & SC-O & Operations and maintenance \\
\midrule
\multirow{4}{*}{\textbf{Technique Integration}}
  & TI-L0 & Single technique (PM only, ML only, or stochastic only) \\
  & TI-L1 & Two techniques present in same study without explicit pipeline contract \\
  & TI-L2 & Partial bridge: two techniques linked, but incomplete or missing ML closure \\
  & TI-L3 & Unified pipeline: PM $\rightarrow$ Markov $\rightarrow$ MC $\rightarrow$ ML
             with formal inter-stage contracts \\
\bottomrule
\end{tabular}
\end{table}

The L0--L3 scale in dimension TI maps to the integration-level coding used
in Sections~\ref{sec:slr-rq3} and~\ref{app:slr-top30}. PATHCAST is positioned
at TI-L3, AD4, with SC coverage spanning SC-D, SC-Q, and SC-C.


% =============================================================================
% APPENDIX E — Reading List and Pipeline Artefacts
% =============================================================================
\chapter{Reading List and Pipeline Artefacts}
\label{app:E}

This appendix presents the ranked reading list of 32 primary studies
selected for in-depth synthesis and documents the scripts and persistent
artefacts of the review pipeline.

\section{Multi-Criteria Scoring Rubric}
\label{app:slr-ranking-criteria}

Table~\ref{tab:app-ranking-criteria} presents the 15-point rubric used to
rank studies. Three seminal papers (Cook 1995, Cook 1998, Rubin 2007) are
included regardless of score.

\begin{table}[htbp]
\centering
\caption{Multi-criteria scoring rubric for the top-32 reading list.}
\label{tab:app-ranking-criteria}
\small
\setlength{\tabcolsep}{4pt}
\begin{tabular}{p{7.5cm} c}
\toprule
\textbf{Criterion} & \textbf{Points} \\
\midrule
IC2 present (stochastic modelling in SE)          & +3 \\
IC3 present (forecasting of process metrics)      & +3 \\
IC1 present (process mining in SE)                & +2 \\
IC4 present (repository mining for process)       & +1 \\
PDF locally available                             & +2 \\
Published in 2020 or later                        & +1 \\
High-impact venue (IEEE/ACM/IST/JSS/IS/TSE)      & +1 \\
Empirical study type (case\_study or experiment)  & +1 \\
Contribution type = prediction or simulation      & +1 \\
\midrule
\textbf{Maximum score}                            & \textbf{15} \\
\bottomrule
\end{tabular}
\end{table}

\section{Top-32 Ranked Reading List}
\label{app:slr-top30}

\begin{table}[p]
\centering
\caption{Top-32 reading list with score, IC cluster, and main finding
  summary. Papers marked $\dagger$ are forced-inclusion seminal works.}
\label{tab:app-top30}
\small
\setlength{\tabcolsep}{3pt}
\begin{tabular}{c c p{3.8cm} c c p{5.7cm}}
\toprule
\textbf{Rank} & \textbf{Pts} & \textbf{Authors (year)} &
  \textbf{ICs} & \textbf{PDF} & \textbf{Main finding (summary)} \\
\midrule
 1 & 14 & Joshi \& Kahani (2024)~\cite{joshi2024}
         & IC2,IC3,IC4 & \checkmark
         & RL/MDP for GitHub PR outcome prediction \\
 2 & 14 & Li et al.\ (2020)~\cite{li2020}
         & IC2,IC3 & \checkmark
         & Hybrid SD+DES simulation for industrial project forecasting \\
 3 & 13 & Guinea-Cabrera et al.\ (2025)~\cite{guinea2025}
         & IC1,IC2,IC3 & \checkmark
         & PRIMAD Digital Twin for SDLC (pm4py + Akka + LightGBM), L2 \\
 4 & 13 & Ortu et al.\ (2023)~\cite{ortu2023}
         & IC2,IC4 & \checkmark
         & Fault-insertion/fixing Markov + SNA in Apache projects \\
 5 & 13 & Ben Mesmia et al.\ (2021)~\cite{benmesmia2021}
         & IC2,IC3 & \checkmark
         & Non-Markovian SPN for DevOps workflow duration prediction \\
 6 & 12 & Buliga et al.\ (2025)~\cite{buliga2025}
         & IC1,IC3 & \checkmark
         & Generative AI + business process simulation for what-if \\
 7 & 12 & Bhadra et al.\ (2023)~\cite{bhadra2023}
         & IC2,IC3 & \checkmark
         & SPN model of CI/CD with reliability sensitivity analysis \\
 8 & 12 & Bhadra (2022)~\cite{bhadra2022}
         & IC2,IC3 & \checkmark
         & SPN model of CI/CD pipeline delivery probability \\
 9 & 12 & Massitela et al.\ (2018)~\cite{massitela2018}
         & IC2,IC3 & \checkmark
         & Stochastic Automata Network for Waterfall team estimation \\
10 & 11 & Incerto et al.\ (2025)~\cite{incerto2025}
         & IC1,IC2 & \checkmark
         & VLMC stochastic conformance; outperforms 18 state-of-art techniques \\
11 & 11 & Jo \& Kwon (2024)~\cite{jo2024}
         & IC2,IC4 & \checkmark
         & Probabilistic model checking (CTL) on GitHub repos \\
12 & 11 & Jo et al.\ (2023)~\cite{jo2023}
         & IC2,IC4 & \checkmark
         & DTMC collaboration patterns in GitHub (5 repos) \\
13 & 11 & L\'opez-Pintado \& Dumas (2023)~\cite{lopezpintado2023}
         & IC1,IC2 & ---
         & Probabilistic resource calendars discovered from event logs \\
14 & 11 & Caldeira et al.\ (2022)~\cite{caldeira2022}
         & IC1,IC3 & \checkmark
         & Process complexity $\to$ cyclomatic complexity ($r \approx 0.43$) \\
15 & 11 & Jalote \& Kamma (2021)~\cite{jalote2021}
         & IC1,IC2 & \checkmark
         & Markov chains on IDE logs for programmer productivity \\
16 & 11 & Lunesu et al.\ (2021)~\cite{lunesu2021}
         & IC2,IC3 & \checkmark
         & Monte Carlo risk assessment for agile software development \\
17 & 11 & Nafreen et al.\ (2020)~\cite{nafreen2020}
         & IC2,IC3 & ---
         & SRGM + defect tracking for hybrid reliability forecasting \\
18 & 11 & Gupta (2017)~\cite{gupta2017}
         & IC1,IC3,IC4 & ---
         & Inductive miner + predictive analytics for SW maintenance \\
19 & 10 & Salmani et al.\ (2022)
         & IC1,IC3 & \checkmark
         & Process mining for requirements engineering in MDPs \\
20 & 10 & Jeon et al.\ (2015)~\cite{jeon2015}
         & IC2,IC3,IC4 & ---
         & Markov-based probabilistic risk prediction from repo data \\
21 & 10 & Tian \& Noore (2005)~\cite{tian2005}
         & IC2,IC3 & ---
         & Markov defect-removal modelling under code churn \\
22 & 10 & Cook \& Wolf (1995)$^\dagger$~\cite{cook1995}
         & IC1,IC2 & \checkmark
         & First automated process discovery from event traces \\
23 &  9 & Pourbafrani et al.\ (2021)~\cite{pourbafrani2021}
         & IC1,IC3 & \checkmark
         & SIMPT: process tree + interactive what-if simulation tool \\
24 &  9 & Tyagi et al.\ (2021)~\cite{tyagi2021}
         & IC2,IC3 & ---
         & State-space Markov framework for agile planning and tracking \\
25 &  9 & Razi et al.\ (2019)
         & IC1,IC2 & \checkmark
         & Conformance checking for SDLC process compliance \\
26 &  9 & Krismayer et al.\ (2019)
         & IC1,IC4 & \checkmark
         & Declarative PM for repository-based process monitoring \\
27 &  9 & Baskoro et al.\ (2019)
         & IC3,IC4 & ---
         & ML-based effort prediction from repository data \\
28 &  9 & (Anon.\ 2017)
         & IC2,IC3 & ---
         & Probabilistic defect prediction for software product lines \\
29 &  9 & Magennis (2015)~\cite{magennis2015}
         & IC2,IC3 & ---
         & Monte Carlo cycle-time economic analysis \\
30 &  9 & Washizaki et al.\ (2015)~\cite{washizaki2015}
         & IC2,IC3 & ---
         & GSRM release-time prediction for open-source projects \\
31 &  9 & Cook \& Wolf (1998)$^\dagger$~\cite{cook1998}
         & IC1,IC2 & ---
         & Markov/FSA/NN process discovery from event-based data \\
32 &  2 & Rubin et al.\ (2007)$^\dagger$~\cite{rubin2007}
         & IC1 & \checkmark
         & First named PM framework for SDLC (ProM + SCM logs) \\
\bottomrule
\end{tabular}
\end{table}

\section{Review Pipeline Scripts}
\label{app:slr-pipeline}

The review was conducted using a custom Python pipeline (\texttt{pipeline/})
operating on a structured CSV data store. Table~\ref{tab:app-pipeline}
documents each script with its function and main invocation command.

\begin{table}[htbp]
\centering
\caption{Review pipeline scripts and their roles.}
\label{tab:app-pipeline}
\small
\setlength{\tabcolsep}{3pt}
\begin{tabular}{p{4.2cm} p{4.5cm} p{4.0cm}}
\toprule
\textbf{Script} & \textbf{Function} & \textbf{Main command} \\
\midrule
\texttt{pipeline/dedup.py} &
  Cross-source deduplication (DOI + fuzzy title) &
  \texttt{python main.py dedup} \\[2pt]
\texttt{pipeline/enrich.py} &
  Abstract enrichment cascade (4 sources) &
  \texttt{python main.py enrich} \\[2pt]
\texttt{pipeline/screening.py} &
  T/A LLM screening via Batches API &
  \texttt{python main.py screen -{}-run} \\[2pt]
\texttt{pipeline/fulltext.py} &
  FT LLM screening + manual review &
  \texttt{python main.py fulltext -{}-run} \\[2pt]
\texttt{pipeline/pdf\_downloader.py} &
  PDF retrieval (Unpaywall, S2, OA) &
  \texttt{python main.py pdf} \\[2pt]
\texttt{pipeline/extract\_prep.py} &
  S2 metadata enrichment + PDF copy &
  \texttt{python pipeline/} \texttt{extract\_prep.py} \\[2pt]
\texttt{pipeline/extract\_llm.py} &
  Structured data extraction (LLM, 169 papers) &
  \texttt{python pipeline/} \texttt{extract\_llm.py -{}-run} \\[2pt]
\texttt{pipeline/synth\_llm.py} &
  Synthesis paragraph generation per IC cluster &
  \texttt{python pipeline/} \texttt{synth\_llm.py} \\[2pt]
\texttt{pipeline/finalization.py} &
  PRISMA snapshot, EC5 closure, artefact export &
  \texttt{python main.py finalize} \\
\bottomrule
\end{tabular}
\end{table}

\section{Persistent Artefact Registry}
\label{app:slr-artefacts}

Table~\ref{tab:app-artefacts} lists the persistent artefacts produced by the
review, with their role, location, and current state.

\begin{table}[htbp]
\centering
\caption{Persistent artefacts of the review.}
\label{tab:app-artefacts}
\small
\setlength{\tabcolsep}{3pt}
\begin{tabular}{p{6.5cm} p{3.0cm} p{2.8cm}}
\toprule
\textbf{File} & \textbf{Role} & \textbf{State} \\
\midrule
\texttt{results/screening/} \texttt{ft\_screening\_results.csv}
  & 886-paper decision log & Final \\[3pt]
\texttt{results/extraction/} \texttt{extraction\_template.csv}
  & 169-paper extraction data & 169/169 \\[3pt]
\texttt{results/final\_review/} \texttt{included\_studies\_current.csv}
  & Included study set & 169 rows \\[3pt]
\texttt{results/final\_review/} \texttt{top30\_reading\_list.csv}
  & Ranked reading list & 32 rows \\[3pt]
\texttt{results/final\_review/} \texttt{cap3\_synthesis\_v2.tex}
  & LLM synthesis paragraphs & For review \\[3pt]
\texttt{results/final\_review/} \texttt{missing\_references.bib}
  & BibTeX for 32 missing keys & Ready to merge \\[3pt]
\texttt{results/final\_review/} \texttt{pending\_inaccessible\_closed.csv}
  & EC5 closure record & 148 rows \\[3pt]
\texttt{results/final\_review/} \texttt{prisma\_summary\_current.txt}
  & PRISMA snapshot & Final \\[3pt]
\texttt{results/} \texttt{SLR\_EXECUCAO\_ATE\_AGORA.md}
  & Execution log (all sessions) & Up to date \\
\bottomrule
\end{tabular}
\end{table}
